% !TeX program = lualatex
\documentclass[12pt,aspectratio=169]{beamer}
\usepackage[utf8]{inputenc}
\usepackage[T1]{fontenc}
\usepackage{amsmath,amsfonts,amssymb}
\usepackage{graphicx}
\usepackage{xcolor}
\usepackage{hyperref}
\usepackage{anyfontsize}
\usepackage{lmodern}
\usepackage{longtable}
\usepackage{array}
\usepackage{booktabs}

% ====== FONT SIZES ======
\renewcommand{\tiny}{\fontsize{4}{5}\selectfont}
\renewcommand{\scriptsize}{\fontsize{5}{6}\selectfont}
\renewcommand{\footnotesize}{\fontsize{6}{7}\selectfont}
\renewcommand{\small}{\fontsize{7}{8}\selectfont}
\renewcommand{\normalsize}{\fontsize{8}{9}\selectfont}
\renewcommand{\large}{\fontsize{9}{10}\selectfont}
\renewcommand{\Large}{\fontsize{10}{11}\selectfont}
\renewcommand{\LARGE}{\fontsize{11}{13}\selectfont}
\renewcommand{\huge}{\fontsize{12}{14}\selectfont}
\renewcommand{\Huge}{\fontsize{14}{17}\selectfont}

% ====== COLORS ======
\definecolor{redcorrect}{RGB}{200,30,30}
\definecolor{bluecorrect}{RGB}{30,80,200}
\definecolor{greencheck}{RGB}{0,150,50}
\definecolor{lightgray}{RGB}{245,245,245}
\definecolor{darkgray}{RGB}{50,50,50}
\definecolor{softyellow}{RGB}{255,248,200}
\definecolor{deepblue}{RGB}{20,60,150}
\definecolor{darkred}{RGB}{150,20,20}
\definecolor{orangewarning}{RGB}{255,140,0}
\definecolor{correctgreen}{RGB}{0,120,40}
\definecolor{trapcolor}{RGB}{180,80,0}

\setbeamercolor{title}{fg=white,bg=redcorrect}
\setbeamercolor{frametitle}{fg=white,bg=redcorrect}
\setbeamercolor{block title}{fg=white,bg=bluecorrect}
\setbeamercolor{block body}{fg=darkgray,bg=lightgray}
\setbeamercolor{normal text}{fg=darkgray}
\usecolortheme{default}

% ====== CUSTOM COMMANDS ======
\newcommand{\correct}[1]{\textcolor{greencheck}{\textbf{\checkmark}} #1}
\newcommand{\keypoint}[1]{\textcolor{deepblue}{\textbf{#1}}}
\newcommand{\answerbox}[1]{\fcolorbox{correctgreen}{green!8}{\parbox{0.88\textwidth}{\centering \textcolor{correctgreen}{\textbf{\checkmark\ CORRECT:}} #1}}}
\newcommand{\trapbox}[1]{\fcolorbox{trapcolor}{orange!8}{\parbox{0.88\textwidth}{\centering \textcolor{trapcolor}{\textbf{⚠ EXAM TRAP:}} #1}}}
\newcommand{\stepbox}[1]{%
	\begin{center}
		\fcolorbox{darkgray}{white}{%
			\parbox{0.90\textwidth}{\vspace{0.05cm} #1 \vspace{0.05cm}}%
		}
	\end{center}
}
\newcommand{\recall}[1]{\fcolorbox{deepblue}{blue!5}{\parbox{0.88\textwidth}{\centering \textbf{REMEMBER:} #1}}}
\newcommand{\wronganswer}[1]{\textcolor{redcorrect}{\textbf{✘}} #1}

% ================================================================
% BEGIN DOCUMENT
% ================================================================
\begin{document}
	
	% ================================================================
	% TITLE SLIDE
	% ================================================================
	\begin{frame}
		\centering
		\vspace{1.2cm}
		{\fontsize{24}{28}\selectfont \textbf{CAMS Domain A}}\\[0.3cm]
		{\fontsize{18}{22}\selectfont \textbf{All PNG Questions Explained}}\\[0.5cm]
		{\fontsize{12}{14}\selectfont \textit{Based on Real Life Questions}}\\[0.3cm]
		\rule{5cm}{0.5pt}\\[0.2cm]
	
	\end{frame}
	
	% ================================================================
	% TABLE OF CONTENTS
	% =================================================== ================================================================
	% SECTION 1: al branch and subsidairie.png
	% ================================================================
	\begin{frame}
		\frametitle{\fontsize{12}{14}\selectfont \textbf{Q1: Recommendation 18 — Group-wide AML Programs}}
		\begin{block}
			\scriptsize
			A multinational institution is criticized during an audit for deficiencies in Recommendation 18 implementation. What should the auditors have likely focused on?
			\vspace{0.1cm}
			\begin{itemize}
				\item[\textbf{A.}] The training curriculum for non-AML business units.
				\item[\textbf{B.}] Whether group-wide AML programs apply to all branches and subsidiaries.
				\item[\textbf{C.}] The marketing strategy for new financial products.
				\item[\textbf{D.}] Customer satisfaction scores for account opening.
			\end{itemize}
		\end{block}
		\vspace{0.05cm}
		\answerbox{\large \textbf{B} is correct.}
		\vspace{0.05cm}
		\stepbox{\scriptsize
			\keypoint{Explanation:} FATF Recommendation 18 requires group-wide AML/CFT programs to apply to \textbf{all branches and subsidiaries}, including those in foreign jurisdictions. Auditors reviewing Recommendation 18 implementation focus on whether the institution's group-wide AML program \textbf{covers all branches and subsidiaries globally}, not just training programs or other operational areas.
		}
		\vspace{0.03cm}
		\trapbox{\scriptsize Rec 18 = \textbf{group-wide AML programs} applying to ALL branches and subsidiaries globally.}
	\end{frame}
	
	% ================================================================
	% SECTION 2: basel aml.png
	% ================================================================
	\begin{frame}
		\frametitle{\fontsize{12}{14}\selectfont \textbf{Q2: Basel AML Index — Vulnerability Indicators}}
		\begin{block}
			\scriptsize
			A national enforcement agency is prioritizing policy reform and wants to identify vulnerable legal structures contributing to money laundering. Why might they choose the Basel AML Index?
			\vspace{0.1cm}
			\begin{itemize}
				\item[\textbf{A.}] It ranks countries by GDP growth rate.
				\item[\textbf{B.}] It compiles vulnerability indicators across legal, political, and financial dimensions.
				\item[\textbf{C.}] It tracks individual money laundering convictions.
				\item[\textbf{D.}] It monitors cryptocurrency exchange volumes.
			\end{itemize}
		\end{block}
		\vspace{0.05cm}
		\answerbox{\large \textbf{B} is correct.}
		\vspace{0.05cm}
		\stepbox{\scriptsize
			\keypoint{Explanation:} The \textbf{Basel AML Index} is a composite index that ranks countries based on their AML/CFT risk. It compiles \textbf{vulnerability indicators across legal, political, and financial dimensions} to provide a comprehensive country risk assessment. It does NOT focus on GDP, convictions, or cryptocurrency alone.
		}
	\end{frame}
	
	% ================================================================
	% Q3: FATF Immediate Outcome 5
	% ================================================================
	\begin{frame}
		\frametitle{\fontsize{12}{14}\selectfont \textbf{Q3: FATF Immediate Outcome 5 — Legal Persons}}
		\begin{block}
			\scriptsize
			A country's high volume of legal arrangements relies on anonymous shell entities. National authorities recently conducted a regulatory sweep but still struggle to trace beneficial ownership. Which FATF Immediate Outcome would be rated as weak?
			\vspace{0.1cm}
			\begin{itemize}
				\item[\textbf{A.}] Immediate Outcome 1 — ML/TF risks are understood.
				\item[\textbf{B.}] Immediate Outcome 3 — Supervision.
				\item[\textbf{C.}] Immediate Outcome 5 — Legal persons and arrangements are not misused.
				\item[\textbf{D.}] Immediate Outcome 8 — Confiscation.
			\end{itemize}
		\end{block}
		\vspace{0.05cm}
		\answerbox{\large \textbf{C} is correct.}
		\vspace{0.05cm}
		\stepbox{\scriptsize
			\keypoint{Explanation:} FATF \textbf{Immediate Outcome 5} specifically assesses \textbf{whether legal persons and arrangements are prevented from being misused for ML/TF}. Anonymous shell entities that prevent beneficial ownership tracing directly indicate a weak IO 5. The inability to trace beneficial ownership despite regulatory sweeps shows that legal structures are being misused.
		}
		\vspace{0.03cm}
		\trapbox{\scriptsize IO 5 = \textbf{legal persons and arrangements not misused}. Anonymous shell entities = weak IO 5.}
	\end{frame}
	
	% ================================================================
	% SECTION 3: device intelligence tools.png
	% ================================================================
	\begin{frame}
		\frametitle{\fontsize{12}{14}\selectfont \textbf{Q4: AI-Based Screening — Inconsistent Rejections}}
		\begin{block}
			\scriptsize
			An internal audit reveals a newly implemented AI-based screening system is rejecting transactions inconsistently across business units. What is the best first step for corrective action?
			\vspace{0.1cm}
			\begin{itemize}
				\item[\textbf{A.}] Immediately adjust the AI model parameters.
				\item[\textbf{B.}] Engage the governance committee to assess model performance.
				\item[\textbf{C.}] Disable the AI system and revert to manual screening.
				\item[\textbf{D.}] Train business units on how to override rejections.
			\end{itemize}
		\end{block}
		\vspace{0.05cm}
		\answerbox{\large \textbf{B} is correct.}
		\vspace{0.05cm}
		\stepbox{\scriptsize
			\keypoint{Explanation:} Inconsistent AI model behavior across business units requires \textbf{governance-level oversight} to assess the model's performance, calibration, and fairness before making technical changes. The governance committee should evaluate whether the model is properly calibrated and whether inconsistencies are due to data differences, model drift, or implementation issues.
		}
	\end{frame}
	
	% ================================================================
	% Q5: Device Intelligence Tools
	% ================================================================
	\begin{frame}
		\frametitle{\fontsize{12}{14}\selectfont \textbf{Q5: Device Intelligence Tools — What They Capture}}
		\begin{block}
			\scriptsize
			Which of the following elements do device intelligence tools capture to enhance detection of online financial crimes? (Select Three.)
			\vspace{0.1cm}
			\begin{itemize}
				\item[\textbf{A.}] IP address and geolocation
				\item[\textbf{B.}] Social media accounts and browsing history
				\item[\textbf{C.}] Operating system and installed plugins
				\item[\textbf{D.}] Browser type and screen resolution
			\end{itemize}
		\end{block}
		\vspace{0.05cm}
		\answerbox{\large \textbf{A, C, D} are correct.}
		\vspace{0.05cm}
		\stepbox{\scriptsize
			\keypoint{Explanation:} Device intelligence tools capture \textbf{device fingerprints}:
			\begin{itemize}
				\item \textbf{IP address and geolocation} — location and network information
				\item \textbf{Operating system and installed plugins} — device configuration
				\item \textbf{Browser type and screen resolution} — browser fingerprinting
			\end{itemize}
			\wronganswer{Social media accounts and browsing history} are NOT captured by device intelligence tools — these would be privacy violations and are not typically collected for device fingerprinting.
		}
		\vspace{0.03cm}
		\trapbox{\scriptsize Device intelligence = device fingerprinting, NOT social media or browsing history.}
	\end{frame}
	
	% ================================================================
	% SECTION 4: ecommerce.png
	% ================================================================
	\begin{frame}
		\frametitle{\fontsize{12}{14}\selectfont \textbf{Q6: E-Commerce ML Mitigation Strategies}}
		\begin{block}
			\scriptsize
			Which of the following strategies help e-commerce platforms mitigate risks related to money laundering? (Select Two.)
			\vspace{0.1cm}
			\begin{itemize}
				\item[\textbf{A.}] Implementing seller-rankings based on transaction volume
				\item[\textbf{B.}] Introducing tiered onboarding based on projected sales activity
				\item[\textbf{C.}] Conducting sanctions screening on all sellers during onboarding
				\item[\textbf{D.}] Requiring all payments be made in cryptocurrency
			\end{itemize}
		\end{block}
		\vspace{0.05cm}
		\answerbox{\large \textbf{B \& C} are correct.}
		\vspace{0.05cm}
		\stepbox{\scriptsize
			\keypoint{Explanation:}
			\begin{itemize}
				\item \textbf{Tiered onboarding based on projected sales activity} — scales CDD to risk level (higher volume = more scrutiny)
				\item \textbf{Sanctions screening on all sellers during onboarding} — catches sanctioned entities before they join the platform
			\end{itemize}
			\wronganswer{Seller-rankings based on transaction volume} could actually \textbf{incentivize ML activity} by encouraging high-volume sellers to launder more money. Cryptocurrency does not resolve ML risk.
		}
	\end{frame}
	
	% ================================================================
	% Q7: Token ML Risks
	% ================================================================
	\begin{frame}
		\frametitle{\fontsize{12}{14}\selectfont \textbf{Q7: Token Use in Money Laundering}}
		\begin{block}
			\scriptsize
			Which of the following are accurate statements regarding the use of tokens in money laundering activities? (Select Two.)
			\vspace{0.1cm}
			\begin{itemize}
				\item[\textbf{A.}] Tokens may be exchanged rapidly on decentralized platforms, obscuring source and destination of funds.
				\item[\textbf{B.}] Tokens may be exchanged rapidly on centralized platforms, obscuring source and destination of funds.
				\item[\textbf{C.}] Token transactions are always fully traceable on all platforms.
				\item[\textbf{D.}] Tokens eliminate all ML risk because they are digital.
			\end{itemize}
		\end{block}
		\vspace{0.05cm}
		\answerbox{\large \textbf{A} is the correct answer (only one correct option shown).}
		\vspace{0.05cm}
		\stepbox{\scriptsize
			\keypoint{Explanation:} \textbf{Decentralized platforms} (DEXs) allow rapid token exchanges \textbf{without KYC}, obscuring source and destination. \textbf{Centralized platforms} require KYC and provide tracing, so they do NOT obscure funds in the same way. Tokens do NOT eliminate ML risk and are not always fully traceable (privacy coins).
		}
		\vspace{0.03cm}
		\trapbox{\scriptsize Decentralized platforms = \textbf{anonymity}. Centralized platforms = \textbf{KYC and traceability}.}
	\end{frame}
	
	% ================================================================
	% SECTION 5: esg.png
	% ================================================================
	\begin{frame}
		\frametitle{\fontsize{12}{14}\selectfont \textbf{Q8: ESG-Aligned Compliance Reporting}}
		\begin{block}
			\scriptsize
			Which of the following reporting efforts would be considered examples of ESG-aligned compliance requirements? (Select Three.)
			\vspace{0.1cm}
			\begin{itemize}
				\item[\textbf{A.}] Monitoring account login patterns for fraud
				\item[\textbf{B.}] Publishing carbon emissions by facility
				\item[\textbf{C.}] Conducting algorithmic risk assessments
				\item[\textbf{D.}] Reporting executive salary data across genders
				\item[\textbf{E.}] Reporting supply chain labor standards
			\end{itemize}
		\end{block}
		\vspace{0.05cm}
		\answerbox{\large \textbf{B, D \& E} are correct.}
		\vspace{0.05cm}
		\stepbox{\scriptsize
			\keypoint{Explanation:} ESG (Environmental, Social, Governance) reporting includes:
			\begin{itemize}
				\item \textbf{Environmental:} Carbon emissions by facility (B)
				\item \textbf{Social:} Executive salary data across genders (D) and supply chain labor standards (E)
			\end{itemize}
			\wronganswer{Monitoring fraud and algorithmic risk assessments} are AML/financial crime controls, not ESG reporting requirements.
		}
	\end{frame}
	
	% ================================================================
	% Q9: UK-EEA Payment Processing
	% ================================================================
	\begin{frame}
		\frametitle{\fontsize{12}{14}\selectfont \textbf{Q9: UK-EEA Payment Processing — AML Framework}}
		\begin{block}
			\scriptsize
			A UK-based financial institution wants to process payments for an online gambling platform headquartered in another EEA country. Which action should be taken before proceeding?
			\vspace{0.1cm}
			\begin{itemize}
				\item[\textbf{A.}] Proceed immediately as EEA countries have harmonized AML rules.
				\item[\textbf{B.}] Conduct enhanced due diligence and evaluate the AML framework of the EEA-based partner.
				\item[\textbf{C.}] Refuse all gambling-related payments under UK policy.
				\item[\textbf{D.}] Only process payments if the gambling platform is licensed in the UK.
			\end{itemize}
		\end{block}
		\vspace{0.05cm}
		\answerbox{\large \textbf{B} is correct.}
		\vspace{0.05cm}
		\stepbox{\scriptsize
			\keypoint{Explanation:} Even within the EEA, AML frameworks and enforcement vary. A UK-based institution must \textbf{conduct enhanced due diligence and evaluate the AML framework of the EEA-based partner} before proceeding. The UK's AML regime requires assessing the AML controls of counterparties, especially for higher-risk sectors like gambling.
		}
	\end{frame}
	
	% ================================================================
	% SECTION 6: external.png
	% ================================================================
	\begin{frame}
		\frametitle{\fontsize{12}{14}\selectfont \textbf{Q10: External Sources for Sanctions Screening}}
		\begin{block}
			\scriptsize
			Which types of lists would be external sources essential for a firm's sanctions screening controls? (Select Two.)
			\vspace{0.1cm}
			\begin{itemize}
				\item[\textbf{A.}] FATF blacklists and UN Security Council lists
				\item[\textbf{B.}] Watchlists developed from the firm's previous investigations
				\item[\textbf{C.}] State Department embargo registries
				\item[\textbf{D.}] Internal customer complaint lists
			\end{itemize}
		\end{block}
		\vspace{0.05cm}
		\answerbox{\large \textbf{A \& C} are correct.}
		\vspace{0.05cm}
		\stepbox{\scriptsize
			\keypoint{Explanation:} \textbf{External sources} for sanctions screening include:
			\begin{itemize}
				\item \textbf{FATF blacklists and UN Security Council lists} — international sanctions and high-risk jurisdictions
				\item \textbf{State Department embargo registries} — US embargoed countries and entities
			\end{itemize}
			\wronganswer{Watchlists from firm's previous investigations} and \wronganswer{internal complaint lists} are \textbf{internal} sources, not external sanctions lists.
		}
	\end{frame}
	
	% ================================================================
	% Q11: Internal vs External Sources
	% ================================================================
	\begin{frame}
		\frametitle{\fontsize{12}{14}\selectfont \textbf{Q11: Internal vs External Data Sources}}
		\begin{block}
			\scriptsize
			Which characteristics describe internal versus external sources used during AFC data extraction? (Select Two.)
			\vspace{0.1cm}
			\begin{itemize}
				\item[\textbf{A.}] Internal sources reflect observed customer behavior or system-held values.
				\item[\textbf{B.}] Internal sources include sanctions and adverse media lists.
				\item[\textbf{C.}] External sources include sanction or adverse media lists from agencies.
				\item[\textbf{D.}] External sources include customer transaction history.
			\end{itemize}
		\end{block}
		\vspace{0.05cm}
		\answerbox{\large \textbf{A \& C} are correct.}
		\vspace{0.05cm}
		\stepbox{\scriptsize
			\keypoint{Explanation:}
			\begin{itemize}
				\item \textbf{Internal sources:} Observed customer behavior, system-held values, transaction history, account data
				\item \textbf{External sources:} Sanction lists, adverse media lists, regulatory databases, PEP lists, watchlists from agencies
			\end{itemize}
		}
	\end{frame}
	
	% ================================================================
	% Q12: Data Lineage
	% ================================================================
	\begin{frame}
		\frametitle{\fontsize{12}{14}\selectfont \textbf{Q12: Data Lineage — Transparency and Auditability}}
		\begin{block}
			\scriptsize
			Which of the following support transparency and auditability through data lineage? (Select Two.)
			\vspace{0.1cm}
			\begin{itemize}
				\item[\textbf{A.}] Archiving analyst notes from alert reviews
				\item[\textbf{B.}] Maintaining history of data transformations
				\item[\textbf{C.}] Capturing source of each data field
				\item[\textbf{D.}] Deleting outdated data annually
			\end{itemize}
		\end{block}
		\vspace{0.05cm}
		\answerbox{\large \textbf{B \& C} are correct.}
		\vspace{0.05cm}
		\stepbox{\scriptsize
			\keypoint{Explanation:} \textbf{Data lineage} supports transparency by:
			\begin{itemize}
				\item \textbf{Maintaining history of data transformations} — shows how data was modified
				\item \textbf{Capturing source of each data field} — shows where data originated
			\end{itemize}
			\wronganswer{Archiving analyst notes} is about case documentation, not data lineage. \wronganswer{Deleting outdated data} undermines auditability.
		}
	\end{frame}
	
	% ================================================================
	% SECTION 7: human traficing topology.png
	% ================================================================
	\begin{frame}
		\frametitle{\fontsize{12}{14}\selectfont \textbf{Q13: Human Trafficking — Financial Indicators}}
		\begin{block}
			\scriptsize
			Authorities investigating a failed terror attack discover that the attackers relied on an informal courier network to move cash across regional borders and stored money temporarily in small hotel safes. What issue does this highlight?
			\vspace{0.1cm}
			\begin{itemize}
				\item[\textbf{A.}] Inadequate cash-flow modeling in foreign exchange settlements
				\item[\textbf{B.}] Vulnerabilities in bulk cash movement and decentralized storage methods
				\item[\textbf{C.}] Weak cryptocurrency regulation
				\item[\textbf{D.}] Inadequate trade finance controls
			\end{itemize}
		\end{block}
		\vspace{0.05cm}
		\answerbox{\large \textbf{B} is correct.}
		\vspace{0.05cm}
		\stepbox{\scriptsize
			\keypoint{Explanation:} Informal courier networks and hotel safe storage highlight \textbf{vulnerabilities in bulk cash movement and decentralized storage methods}. These methods bypass the formal financial system entirely, avoiding AML/CFT controls, transaction monitoring, and reporting thresholds.
		}
	\end{frame}
	
	% ================================================================
	% Q14: Human Trafficking Risks
	% ================================================================
	\begin{frame}
		\frametitle{\fontsize{12}{14}\selectfont \textbf{Q14: Human Trafficking AML Risks}}
		\begin{block}
			\scriptsize
			Which of the following are recognized risks associated with human trafficking and human smuggling in the context of AML efforts? (Select Three.)
			\vspace{0.1cm}
			\begin{itemize}
				\item[\textbf{A.}] Use of trafficked individuals for value transfer
				\item[\textbf{B.}] Variability in how jurisdictions classify exploitation
				\item[\textbf{C.}] Decrease in reporting due to legitimate migration protections
				\item[\textbf{D.}] Misclassification under predicate offenses in FATF standards
			\end{itemize}
		\end{block}
		\vspace{0.05cm}
		\answerbox{\large \textbf{A, B \& D} are correct.}
		\vspace{0.05cm}
		\stepbox{\scriptsize
			\keypoint{Explanation:}
			\begin{itemize}
				\item \textbf{A:} Trafficked individuals may be used as mules to transfer value
				\item \textbf{B:} Jurisdictions vary in how they classify exploitation (labor vs sexual vs other)
				\item \textbf{D:} Misclassification under predicate offenses makes it harder to identify trafficking
			\end{itemize}
			\wronganswer{Decrease in reporting due to migration protections} is NOT a recognized risk — migration protections would increase reporting and protection for victims.
		}
	\end{frame}
	
	% ================================================================
	% SECTION 8: icar.png
	% ================================================================
	\begin{frame}
		\frametitle{\fontsize{12}{14}\selectfont \textbf{Q15: ICAR Legal Advisory Services}}
		\begin{block}
			\scriptsize
			An NGO is developing anti-corruption mechanisms in a country with weak legal infrastructure and endemic fraud. To effectively implement reforms, which combination of Basel Institute services would provide the most strategic support?
			\vspace{0.1cm}
			\begin{itemize}
				\item[\textbf{A.}] Public finance management technical assistance with standardized compliance templates
				\item[\textbf{B.}] ICAR legal advisory services combined with capacity building programs
				\item[\textbf{C.}] Third-party auditing services only
				\item[\textbf{D.}] Technology solutions only
			\end{itemize}
		\end{block}
		\vspace{0.05cm}
		\answerbox{\large \textbf{B} is correct.}
		\vspace{0.05cm}
		\stepbox{\scriptsize
			\keypoint{Explanation:} In countries with weak legal infrastructure, \textbf{ICAR (International Centre for Asset Recovery) legal advisory services combined with capacity building programs} provide the most strategic support. Legal advisory addresses the framework gaps, while capacity building ensures sustainable implementation. Standardized templates or auditing alone would not address the underlying legal weaknesses.
		}
	\end{frame}
	
	% ================================================================
	% SECTION 9: iosco.png
	% ================================================================
	\begin{frame}
		\frametitle{\fontsize{12}{14}\selectfont \textbf{Q16: IOSCO Core Objectives}}
		\begin{block}
			\scriptsize
			Which of IOSCO's core objectives directly support anti-financial crime efforts in securities markets? (Select Two.)
			\vspace{0.1cm}
			\begin{itemize}
				\item[\textbf{A.}] Enhancing investor protection through cooperation and enforcement
				\item[\textbf{B.}] Requiring all securities firms to maintain compliance departments with uniform structures
				\item[\textbf{C.}] Promoting financial stability by managing systemic risks and facilitating international information exchange
				\item[\textbf{D.}] Standardizing all securities products globally
			\end{itemize}
		\end{block}
		\vspace{0.05cm}
		\answerbox{\large \textbf{A \& C} are correct.}
		\vspace{0.05cm}
		\stepbox{\scriptsize
			\keypoint{Explanation:} IOSCO's core objectives supporting AFC:
			\begin{itemize}
				\item \textbf{Enhancing investor protection through cooperation and enforcement} — helps combat market abuse and fraud
				\item \textbf{Promoting financial stability by managing systemic risks and facilitating international information exchange} — supports cross-border financial crime detection
			\end{itemize}
			\wronganswer{Uniform compliance departments} and \wronganswer{global product standardization} are NOT IOSCO core objectives.
		}
	\end{frame}
	
	% ================================================================
	% SECTION 10: open source research.png
	% ================================================================
	\begin{frame}
		\frametitle{\fontsize{12}{14}\selectfont \textbf{Q17: Open-Source Research for Investigations}}
		\begin{block}
			\scriptsize
			While reviewing an alert, the analyst notes wire transfers to an unfamiliar entity in a sanctioned jurisdiction. The customer profile does not list dealings in that region. Which information source would most help the analyst understand the activity?
			\vspace{0.1cm}
			\begin{itemize}
				\item[\textbf{A.}] Regulatory advisories about common evasion techniques in that jurisdiction
				\item[\textbf{B.}] Open-source research into the entity
				\item[\textbf{C.}] Internal audit reports from the previous year
				\item[\textbf{D.}] General economic data about the jurisdiction
			\end{itemize}
		\end{block}
		\vspace{0.05cm}
		\answerbox{\large \textbf{B} is correct.}
		\vspace{0.05cm}
		\stepbox{\scriptsize
			\keypoint{Explanation:} When encountering an \textbf{unfamiliar entity in a sanctioned jurisdiction}, \textbf{open-source research} into the entity provides the most direct and actionable information. This can reveal the entity's ownership, business activities, connections, and whether it is a shell company or front. Regulatory advisories are helpful but entity-specific research is the most useful next step.
		}
	\end{frame}
	
	% ================================================================
	% Q18: Manual Escalation of Suspicious Activity
	% ================================================================
	\begin{frame}
		\frametitle{\fontsize{12}{14}\selectfont \textbf{Q18: Manual Escalation — Frontline Staff}}
		\begin{block}
			\scriptsize
			A frontline staff receives a large, unusually timed cash deposit from a long-standing customer at a small bank branch. There is no immediate system alert. What is the most appropriate action the staff should take?
			\vspace{0.1cm}
			\begin{itemize}
				\item[\textbf{A.}] Ignore the deposit if no system alert was generated.
				\item[\textbf{B.}] Manually escalate a report to compliance.
				\item[\textbf{C.}] Process the deposit and wait for the next system update.
				\item[\textbf{D.}] Ask the customer to explain the deposit in writing.
			\end{itemize}
		\end{block}
		\vspace{0.05cm}
		\answerbox{\large \textbf{B} is correct.}
		\vspace{0.05cm}
		\stepbox{\scriptsize
			\keypoint{Explanation:} Frontline staff are the \textbf{first line of defense} and must escalate suspicious activity even when there is no system alert. A large, unusually timed cash deposit from a long-standing customer is a red flag that requires \textbf{manual escalation to compliance}. Staff should not ignore, wait, or question the customer directly (tipping off risk).
		}
	\end{frame}
	
	% ================================================================
	% SECTION 11: otc derivative.png
	% ================================================================
	\begin{frame}
		\frametitle{\fontsize{12}{14}\selectfont \textbf{Q19: Post-Trade Surveillance for OTC Derivatives}}
		\begin{block}
			\scriptsize
			A capital markets firm's compliance team identifies gaps in post-trade surveillance for over-the-counter (OTC) derivatives. Which of the following actions would best improve transaction transparency? (Select Two.)
			\vspace{0.1cm}
			\begin{itemize}
				\item[\textbf{A.}] Require traders to document the business purpose for each OTC transaction.
				\item[\textbf{B.}] Implement end-of-day reporting protocols.
				\item[\textbf{C.}] Enable automated reconciliation between front and back office systems.
				\item[\textbf{D.}] Eliminate all manual trade reviews.
			\end{itemize}
		\end{block}
		\vspace{0.05cm}
		\answerbox{\large \textbf{B \& C} are correct.}
		\vspace{0.05cm}
		\stepbox{\scriptsize
			\keypoint{Explanation:}
			\begin{itemize}
				\item \textbf{End-of-day reporting protocols} — ensures all trades are captured and reported
				\item \textbf{Automated reconciliation between front and back office} — identifies discrepancies and ensures transparency
			\end{itemize}
			\wronganswer{Documenting business purpose} is helpful but not structural transparency. \wronganswer{Eliminating manual reviews} could miss important context.
		}
	\end{frame}
	
	% ================================================================
	% SECTION 12: pep.png
	% ================================================================
	\begin{frame}
		\frametitle{\fontsize{12}{14}\selectfont \textbf{Q20: PEP Risk Management Measures}}
		\begin{block}
			\scriptsize
			Which of the following are effective measures a financial institution can implement to manage its exposure to PEP-related money laundering risks? (Select Three.)
			\vspace{0.1cm}
			\begin{itemize}
				\item[\textbf{A.}] Capture information about source of funds and source of wealth at onboarding
				\item[\textbf{B.}] Conduct ongoing media monitoring for adverse news on clients
				\item[\textbf{C.}] Implement processes to adjust transaction thresholds annually for all clients
				\item[\textbf{D.}] Establish a system that auto-detects public official occupations
			\end{itemize}
		\end{block}
		\vspace{0.05cm}
		\answerbox{\large \textbf{A, B \& D} are correct.}
		\vspace{0.05cm}
		\stepbox{\scriptsize
			\keypoint{Explanation:}
			\begin{itemize}
				\item \textbf{A:} Source of funds/wealth is required PEP EDD (FATF Rec 12)
				\item \textbf{B:} Ongoing media monitoring detects adverse news and changes in status
				\item \textbf{D:} Auto-detection of public official occupations identifies PEP status at onboarding
			\end{itemize}
			\wronganswer{Adjusting transaction thresholds annually for all clients} is NOT PEP-specific and could be applied to all clients, not targeting PEP risk specifically.
		}
	\end{frame}
	
	% ================================================================
	% Q21: Prepaid Cards — ML Vulnerability
	% ================================================================
	\begin{frame}
		\frametitle{\fontsize{12}{14}\selectfont \textbf{Q21: Prepaid Cards — ML Vulnerability}}
		\begin{block}
			\scriptsize
			Why are prepaid cards particularly vulnerable to financial crime risks?
			\vspace{0.1cm}
			\begin{itemize}
				\item[\textbf{A.}] They always require full KYC, making them high-risk.
				\item[\textbf{B.}] They can be used anonymously, especially when reloadable or unregistered.
				\item[\textbf{C.}] They have high transaction limits.
				\item[\textbf{D.}] They are not covered by AML regulations.
			\end{itemize}
		\end{block}
		\vspace{0.05cm}
		\answerbox{\large \textbf{B} is correct.}
		\vspace{0.05cm}
		\stepbox{\scriptsize
			\keypoint{Explanation:} Prepaid cards are vulnerable because they can be used \textbf{anonymously, especially when reloadable or unregistered}. This makes them attractive for placement and layering of illicit funds. They are \textbf{covered by AML regulations} in most jurisdictions, but the anonymity feature creates inherent risk.
		}
	\end{frame}
	
	% ================================================================
	% SECTION 13: private banking.png
	% ================================================================
	\begin{frame}
		\frametitle{\fontsize{12}{14}\selectfont \textbf{Q22: Private Banking ML Risk Factors}}
		\begin{block}
			\scriptsize
			Which factor increases the money laundering risk in private banking?
			\vspace{0.1cm}
			\begin{itemize}
				\item[\textbf{A.}] High turnover of low-value retail clients.
				\item[\textbf{B.}] Offshore jurisdictions used to establish client accounts.
				\item[\textbf{C.}] Standardized products for mass market clients.
				\item[\textbf{D.}] Use of domestic payment systems only.
			\end{itemize}
		\end{block}
		\vspace{0.05cm}
		\answerbox{\large \textbf{B} is correct.}
		\vspace{0.05cm}
		\stepbox{\scriptsize
			\keypoint{Explanation:} \textbf{Offshore jurisdictions} used to establish client accounts increase ML risk due to:
			\begin{itemize}
				\item Reduced transparency
				\item Weaker AML enforcement in some jurisdictions
				\item Complexity of cross-border structures
				\item Difficulty in tracing beneficial ownership
			\end{itemize}
		}
	\end{frame}
	
	% ================================================================
	% Q23: High-Risk Private Banking Features
	% ================================================================
	\begin{frame}
		\frametitle{\fontsize{12}{14}\selectfont \textbf{Q23: High-Risk Private Banking Features}}
		\begin{block}
			\scriptsize
			Which of the following could be high-risk features of private banking products or customers? (Select Three.)
			\vspace{0.1cm}
			\begin{itemize}
				\item[\textbf{A.}] Minimum assets requirement
				\item[\textbf{B.}] Offering numbered accounts
				\item[\textbf{C.}] Dedicated relationship service with transaction discretion
				\item[\textbf{D.}] Standardized online account opening
			\end{itemize}
		\end{block}
		\vspace{0.05cm}
		\answerbox{\large \textbf{A, B \& C} are correct.}
		\vspace{0.05cm}
		\stepbox{\scriptsize
			\keypoint{Explanation:}
			\begin{itemize}
				\item \textbf{A:} Minimum assets attract high-net-worth individuals with potential complex financial structures
				\item \textbf{B:} Numbered accounts reduce transparency and conceal identity
				\item \textbf{C:} Dedicated relationship managers with transaction discretion enable circumvention of controls
			\end{itemize}
			\wronganswer{Standardized online account opening} is lower risk, not high-risk.
		}
	\end{frame}
	
	% ================================================================
	% SECTION 14: private banking 2.png
	% ================================================================
	\begin{frame}
		\frametitle{\fontsize{12}{14}\selectfont \textbf{Q24: Synthetic Identity Fraud — Corrective Action}}
		\begin{block}
			\scriptsize
			A new credit product is rolled out targeting low-income customers. Shortly after launch, the firm identifies a trend of synthetic identity fraud linked to this product. Which of the following is the most appropriate corrective action?
			\vspace{0.1cm}
			\begin{itemize}
				\item[\textbf{A.}] Continue the product while conducting a parallel investigation.
				\item[\textbf{B.}] Suspend the product while updating the risk assessment and onboarding controls.
				\item[\textbf{C.}] Increase the credit limit to attract legitimate customers.
				\item[\textbf{D.}] Only investigate individual cases without changing controls.
			\end{itemize}
		\end{block}
		\vspace{0.05cm}
		\answerbox{\large \textbf{B} is correct.}
		\vspace{0.05cm}
		\stepbox{\scriptsize
			\keypoint{Explanation:} When a product is linked to synthetic identity fraud, the most appropriate action is to \textbf{suspend the product while updating the risk assessment and onboarding controls}. This stops further fraud while controls are strengthened. Continuing the product risks additional fraud losses and regulatory exposure.
		}
	\end{frame}
	
	% ================================================================
	% SECTION 15: tech and data.png
	% ================================================================
	\begin{frame}
		\frametitle{\fontsize{12}{14}\selectfont \textbf{Q25: Transitioning to AI-Based Compliance Tools}}
		\begin{block}
			\scriptsize
			Which of the following actions would be most important when transitioning from traditional to AI-based compliance tools to maintain effectiveness and regulatory compliance? (Select Three.)
			\vspace{0.1cm}
			\begin{itemize}
				\item[\textbf{A.}] Verifying transparency and explainability in AI models
				\item[\textbf{B.}] Ensuring data quality for training algorithms
				\item[\textbf{C.}] Selecting third-party global audit providers
				\item[\textbf{D.}] Managing the change in employee workflows
			\end{itemize}
		\end{block}
		\vspace{0.05cm}
		\answerbox{\large \textbf{A, B \& D} are correct.}
		\vspace{0.05cm}
		\stepbox{\scriptsize
			\keypoint{Explanation:}
			\begin{itemize}
				\item \textbf{A:} Transparency/explainability is critical for regulatory compliance
				\item \textbf{B:} Data quality determines AI model effectiveness
				\item \textbf{D:} Managing employee workflow changes ensures adoption and proper use
			\end{itemize}
			\wronganswer{Selecting third-party audit providers} is less critical than the core transition elements — audit providers should be selected during implementation, not as a primary transition action.
		}
	\end{frame}
	
	% ================================================================
	% Q26: RPA in AML Operations
	% ================================================================
	\begin{frame}
		\frametitle{\fontsize{12}{14}\selectfont \textbf{Q26: Robotic Process Automation (RPA) Benefits}}
		\begin{block}
			\scriptsize
			One advantage of robotic process automation (RPA) in AML operations is that it can:
			\vspace{0.1cm}
			\begin{itemize}
				\item[\textbf{A.}] Replace all human decision-making in AML.
				\item[\textbf{B.}] Reduce compliance costs by handling large-volume, repeatable tasks.
				\item[\textbf{C.}] Eliminate all false positives.
				\item[\textbf{D.}] Automate regulatory reporting to all jurisdictions.
			\end{itemize}
		\end{block}
		\vspace{0.05cm}
		\answerbox{\large \textbf{B} is correct.}
		\vspace{0.05cm}
		\stepbox{\scriptsize
			\keypoint{Explanation:} RPA \textbf{reduces compliance costs by handling large-volume, repeatable tasks} (data extraction, reconciliation, report generation). It does \textbf{NOT} replace human decision-making, eliminate false positives, or automate all regulatory reporting across jurisdictions.
		}
	\end{frame}
	
	% ================================================================
	% SECTION 16: six step framework.png
	% ================================================================
	\begin{frame}
		\frametitle{\fontsize{12}{14}\selectfont \textbf{Q27: FATF Six-Step Risk Assessment Framework}}
		\begin{block}
			\scriptsize
			Which of the following are key components of the six-step framework promoted by FATF for conducting national money laundering and terrorist financing risk assessments? (Select Three.)
			\vspace{0.1cm}
			\begin{itemize}
				\item[\textbf{A.}] Environmental scanning
				\item[\textbf{B.}] Horizon scanning
				\item[\textbf{C.}] Threat identification
				\item[\textbf{D.}] Cross-border tax enforcement
			\end{itemize}
		\end{block}
		\vspace{0.05cm}
		\answerbox{\large \textbf{A, B \& C} are correct.}
		\vspace{0.05cm}
		\stepbox{\scriptsize
			\keypoint{Explanation:} The FATF six-step framework for national ML/TF risk assessments includes:
			\begin{itemize}
				\item \textbf{Environmental scanning} — understanding the context
				\item \textbf{Horizon scanning} — identifying emerging risks
				\item \textbf{Threat identification} — identifying specific threats
			\end{itemize}
			\wronganswer{Cross-border tax enforcement} is NOT a component of the FATF risk assessment framework.
		}
	\end{frame}
	
	% ================================================================
	% Q28: Financial Secrecy Index Factors
	% ================================================================
	\begin{frame}
		\frametitle{\fontsize{12}{14}\selectfont \textbf{Q28: Financial Secrecy Index — Elevated Rating}}
		\begin{block}
			\scriptsize
			Which of the following might contribute to a jurisdiction's elevated rating according to the Financial Secrecy Index? (Select Three.)
			\vspace{0.1cm}
			\begin{itemize}
				\item[\textbf{A.}] Minimal reporting of beneficial ownership
				\item[\textbf{B.}] Extensive cross-border services to non-residents
				\item[\textbf{C.}] High ratio of secrecy score to global scale weight
				\item[\textbf{D.}] High ranking on global GDP without AML laws
			\end{itemize}
		\end{block}
		\vspace{0.05cm}
		\answerbox{\large \textbf{A, B \& C} are correct.}
		\vspace{0.05cm}
		\stepbox{\scriptsize
			\keypoint{Explanation:}
			\begin{itemize}
				\item \textbf{A:} Minimal BO reporting increases secrecy
				\item \textbf{B:} Cross-border services to non-residents increase secrecy risk
				\item \textbf{C:} High secrecy score relative to scale weight indicates concentration of secrecy
			\end{itemize}
			\wronganswer{High GDP without AML laws} is not a direct Financial Secrecy Index factor — GDP alone doesn't indicate secrecy.
		}
	\end{frame}
	
	% ================================================================
	% SECTION 17: regulator reporting.png
	% ================================================================
	\begin{frame}
		\frametitle{\fontsize{12}{14}\selectfont \textbf{Q29: Corporate Client EDD — Multiple Layers}}
		\begin{block}
			\scriptsize
			A financial institution is onboarding a new corporate client. During review, the onboarding team identifies that the company's ownership structure includes multiple layers across different jurisdictions, and one of the shareholders is a nominee company based in a recognized secrecy haven. What should be the onboarding team's next step?
			\vspace{0.1cm}
			\begin{itemize}
				\item[\textbf{A.}] Proceed with standard CDD as the structure is legal.
				\item[\textbf{B.}] Escalate the case for enhanced due diligence and verify the identity of the ultimate beneficial owner.
				\item[\textbf{C.}] Decline the client immediately.
				\item[\textbf{D.}] Request the client to simplify the structure.
			\end{itemize}
		\end{block}
		\vspace{0.05cm}
		\answerbox{\large \textbf{B} is correct.}
		\vspace{0.05cm}
		\stepbox{\scriptsize
			\keypoint{Explanation:} Multiple layers, cross-border jurisdictions, and a nominee company in a secrecy haven are \textbf{high-risk indicators} requiring \textbf{escalation for EDD}. The key is to \textbf{verify the identity of the ultimate beneficial owner}. Immediate decline without investigation is not the appropriate FATF risk-based approach.
		}
	\end{frame}
	
	% ================================================================
	% Q30: AFC Regulatory Reporting Practices
	% ================================================================
	\begin{frame}
		\frametitle{\fontsize{12}{14}\selectfont \textbf{Q30: AFC Regulatory Reporting Practices}}
		\begin{block}
			\scriptsize
			Which of the following practices best support accurate and consistent AFC regulatory reporting? (Select Two.)
			\vspace{0.1cm}
			\begin{itemize}
				\item[\textbf{A.}] Using customer segmentation only based on industry codes
				\item[\textbf{B.}] Implementing a centralized transaction monitoring system
				\item[\textbf{C.}] Submitting risk assessments upon regulatory request
				\item[\textbf{D.}] Regular data cleansing of KYC and payment data sources
			\end{itemize}
		\end{block}
		\vspace{0.05cm}
		\answerbox{\large \textbf{B \& D} are correct.}
		\vspace{0.05cm}
		\stepbox{\scriptsize
			\keypoint{Explanation:}
			\begin{itemize}
				\item \textbf{B:} Centralized transaction monitoring ensures consistency
				\item \textbf{D:} Regular data cleansing ensures accuracy of reporting
			\end{itemize}
			\wronganswer{Segmentation based only on industry codes} is insufficient. \wronganswer{Submitting risk assessments only upon request} is reactive, not supportive of consistent reporting.
		}
	\end{frame}
	
	% ================================================================
	% SECTION 18: hnwi.png
	% ================================================================
	\begin{frame}
		\frametitle{\fontsize{12}{14}\selectfont \textbf{Q31: High-Net-Worth Client — Data Review}}
		\begin{block}
			\scriptsize
			A compliance officer finds that transactions of an internal high-net-worth client do not match their profile data (collected via KYC). Which data should they review next to assess risk accurately? (Select Two.)
			\vspace{0.1cm}
			\begin{itemize}
				\item[\textbf{A.}] Current sanctions database
				\item[\textbf{B.}] Counterparty information
				\item[\textbf{C.}] Transaction velocity metrics
				\item[\textbf{D.}] Local news reports
			\end{itemize}
		\end{block}
		\vspace{0.05cm}
		\answerbox{\large \textbf{B \& C} are correct.}
		\vspace{0.05cm}
		\stepbox{\scriptsize
			\keypoint{Explanation:} When transactions don't match the KYC profile:
			\begin{itemize}
				\item \textbf{Counterparty information} — who is the client transacting with?
				\item \textbf{Transaction velocity metrics} — frequency and speed of transactions
			\end{itemize}
			\wronganswer{Sanctions database} is an initial screening step, not the next review for profile mismatches. \wronganswer{Local news reports} are less direct than counterparty/velocity data.
		}
	\end{frame}
	
	% ================================================================
	% SECTION 19: ewra.png
	% ================================================================
	\begin{frame}
		\frametitle{\fontsize{12}{14}\selectfont \textbf{Q32: EWRA Results — MLRO Role}}
		\begin{block}
			\scriptsize
			Which function within an institution uses EWRA results to shape risk posture?
			\vspace{0.1cm}
			\begin{itemize}
				\item[\textbf{A.}] Internal Audit Lead
				\item[\textbf{B.}] Money Laundering Reporting Officer (MLRO)
				\item[\textbf{C.}] Marketing Department
				\item[\textbf{D.}] Human Resources
			\end{itemize}
		\end{block}
		\vspace{0.05cm}
		\answerbox{\large \textbf{B} is correct.}
		\vspace{0.05cm}
		\stepbox{\scriptsize
			\keypoint{Explanation:} The \textbf{MLRO} is responsible for shaping the institution's risk posture based on \textbf{Enterprise-Wide Risk Assessment (EWRA)} results. The MLRO uses EWRA findings to adjust AML/CFT controls, enhance monitoring, and allocate resources to higher-risk areas. Internal Audit assesses controls but doesn't shape risk posture.
		}
	\end{frame}
	
	% ================================================================
	% Q33: Financial Crime Risk Profile
	% ================================================================
	\begin{frame}
		\frametitle{\fontsize{12}{14}\selectfont \textbf{Q33: Financial Crime Risk Profile — High-Risk Factors}}
		\begin{block}
			\scriptsize
			Which of the following conditions increase a financial institution's financial crime risk profile? (Select Two.)
			\vspace{0.1cm}
			\begin{itemize}
				\item[\textbf{A.}] Operating in regions with weak regulatory enforcement
				\item[\textbf{B.}] Offering offshore private banking services
				\item[\textbf{C.}] Having a strong compliance culture
				\item[\textbf{D.}] Maintaining low customer onboarding thresholds
			\end{itemize}
		\end{block}
		\vspace{0.05cm}
		\answerbox{\large \textbf{A \& B} are correct.}
		\vspace{0.05cm}
		\stepbox{\scriptsize
			\keypoint{Explanation:}
			\begin{itemize}
				\item \textbf{A:} Weak regulatory enforcement increases inherent risk
				\item \textbf{B:} Offshore private banking services increase ML risk
			\end{itemize}
			\wronganswer{Strong compliance culture} reduces risk, doesn't increase it. \wronganswer{Low onboarding thresholds} may actually reduce risk by enabling better monitoring.
		}
	\end{frame}
	
	% ================================================================
	% Q34: Tax Evasion Conclusion
	% ================================================================
	\begin{frame}
		\frametitle{\fontsize{12}{14}\selectfont \textbf{Q34: Tax Evasion vs Legal Avoidance}}
		\begin{block}
			\scriptsize
			An internal audit reveals that a company failed to report offshore revenue and laundered funds through an unrelated third party. Which of the following conclusions is best supported by this evidence?
			\vspace{0.1cm}
			\begin{itemize}
				\item[\textbf{A.}] The company legally avoided taxes through offshore structuring.
				\item[\textbf{B.}] The company committed tax evasion by concealing income.
				\item[\textbf{C.}] The company engaged in legitimate tax planning.
				\item[\textbf{D.}] The company's actions are legal under international tax law.
			\end{itemize}
		\end{block}
		\vspace{0.05cm}
		\answerbox{\large \textbf{B} is correct.}
		\vspace{0.05cm}
		\stepbox{\scriptsize
			\keypoint{Explanation:} \textbf{Failure to report offshore revenue + laundering through a third party} = \textbf{tax evasion by concealing income}. This is NOT legal tax avoidance — tax avoidance requires full disclosure. Concealment plus laundering indicates criminal tax evasion, a predicate offense for ML.
		}
	\end{frame}
	
	% ================================================================
	% SECTION 20: crypto.png
	% ================================================================
	\begin{frame}
		\frametitle{\fontsize{12}{14}\selectfont \textbf{Q35: Payment Screening — Hybrid System}}
		\begin{block}
			\scriptsize
			A global bank with a high transaction volume is updating its payment screening system. Its goal is to minimize false positives while maintaining regulatory compliance. Which configuration of tools would best achieve this?
			\vspace{0.1cm}
			\begin{itemize}
				\item[\textbf{A.}] Use a single rule-based filter for all payments.
				\item[\textbf{B.}] Use a hybrid system combining fuzzy logic, behavioral AI models, and rule-based filters.
				\item[\textbf{C.}] Use AI only, removing all rules.
				\item[\textbf{D.}] Use manual review for all payments.
			\end{itemize}
		\end{block}
		\vspace{0.05cm}
		\answerbox{\large \textbf{B} is correct.}
		\vspace{0.05cm}
		\stepbox{\scriptsize
			\keypoint{Explanation:} A \textbf{hybrid system combining fuzzy logic, behavioral AI models, and rule-based filters} best balances false positives and compliance. Fuzzy logic handles partial matches, behavioral AI detects anomalies, and rules provide regulatory coverage. Single approaches (rules only, AI only, or manual) are less effective.
		}
	\end{frame}
	
	% ================================================================
	% Q36: Blockchain Transaction Tracing
	% ================================================================
	\begin{frame}
		\frametitle{\fontsize{12}{14}\selectfont \textbf{Q36: Blockchain Transaction Tracing Methods}}
		\begin{block}
			\scriptsize
			Which of the following technologies or methods support blockchain transaction tracing activities? (Select Two.)
			\vspace{0.1cm}
			\begin{itemize}
				\item[\textbf{A.}] Distributed IPFS file sharing
				\item[\textbf{B.}] Wallet clustering algorithms
				\item[\textbf{C.}] Chain visualization dashboards
				\item[\textbf{D.}] Centralized exchange order books
			\end{itemize}
		\end{block}
		\vspace{0.05cm}
		\answerbox{\large \textbf{B \& C} are correct.}
		\vspace{0.05cm}
		\stepbox{\scriptsize
			\keypoint{Explanation:}
			\begin{itemize}
				\item \textbf{Wallet clustering algorithms} — group related wallets to identify entities
				\item \textbf{Chain visualization dashboards} — visually trace transaction flows
			\end{itemize}
			\wronganswer{IPFS} is a file system, not tracing. \wronganswer{Order books} are exchange-specific, not blockchain tracing.
		}
	\end{frame}
	
	% ================================================================
	% SECTION 21: bsa.png
	% ================================================================
	\begin{frame}
		\frametitle{\fontsize{12}{14}\selectfont \textbf{Q37: BSA AML/CFT Policies — Key Features}}
		\begin{block}
			\scriptsize
			Which of the following are key features of well-structured AML/CFT policies and procedures aligned with the Bank Secrecy Act? (Select Two.)
			\vspace{0.1cm}
			\begin{itemize}
				\item[\textbf{A.}] Requiring transaction reporting above a fixed institutional threshold
				\item[\textbf{B.}] Tailoring procedures to specific business units or jurisdictions
				\item[\textbf{C.}] Assigning roles and responsibilities to specific individuals or teams
				\item[\textbf{D.}] Using the same procedures for all business units globally
			\end{itemize}
		\end{block}
		\vspace{0.05cm}
		\answerbox{\large \textbf{B \& C} are correct.}
		\vspace{0.05cm}
		\stepbox{\scriptsize
			\keypoint{Explanation:}
			\begin{itemize}
				\item \textbf{B:} Tailoring procedures to specific business units or jurisdictions (risk-based approach)
				\item \textbf{C:} Assigning roles and responsibilities to specific individuals or teams (accountability)
			\end{itemize}
			\wronganswer{Fixed thresholds} are not risk-based. \wronganswer{Same procedures globally} ignores jurisdictional risk differences.
		}
	\end{frame}
	
	% ================================================================
	% Q38: China AML Law Evolution
	% ================================================================
	\begin{frame}
		\frametitle{\fontsize{12}{14}\selectfont \textbf{Q38: China AML Law 2025 — Expanded Coverage}}
		\begin{block}
			\scriptsize
			What is an accurate statement regarding China's evolution of AML Law which took effect in 2025?
			\vspace{0.1cm}
			\begin{itemize}
				\item[\textbf{A.}] It reduced coverage to only banks.
				\item[\textbf{B.}] It added coverage for additional sectors including law firms, real estate agencies, and dealers in precious gems.
				\item[\textbf{C.}] It eliminated all reporting requirements.
				\item[\textbf{D.}] It only applies to state-owned enterprises.
			\end{itemize}
		\end{block}
		\vspace{0.05cm}
		\answerbox{\large \textbf{B} is correct.}
		\vspace{0.05cm}
		\stepbox{\scriptsize
			\keypoint{Explanation:} China's 2025 AML Law \textbf{expanded coverage} to include \textbf{designated non-financial businesses and professions (DNFBPs)} including law firms, real estate agencies, and dealers in precious gems and metals. This aligns China with FATF standards requiring DNFBPs to implement AML/CFT obligations.
		}
	\end{frame}
	
	% ================================================================
	% SECTION 22: clustering.png
	% ================================================================
	\begin{frame}
		\frametitle{\fontsize{12}{14}\selectfont \textbf{Q39: Clustering in Transaction Monitoring}}
		\begin{block}
			\scriptsize
			What is the primary benefit of clustering within transaction monitoring systems?
			\vspace{0.1cm}
			\begin{itemize}
				\item[\textbf{A.}] It reduces all false positives to zero.
				\item[\textbf{B.}] It reveals shared patterns among entities to detect potential financial crime.
				\item[\textbf{C.}] It automatically files SARs.
				\item[\textbf{D.}] It eliminates the need for manual review.
			\end{itemize}
		\end{block}
		\vspace{0.05cm}
		\answerbox{\large \textbf{B} is correct.}
		\vspace{0.05cm}
		\stepbox{\scriptsize
			\keypoint{Explanation:} \textbf{Clustering} reveals \textbf{shared patterns among entities} — linking otherwise separate accounts or transactions that exhibit similar behaviors. This helps detect networks of criminal activity (e.g., money mule networks, organized crime). It does NOT eliminate false positives, file SARs automatically, or replace manual review.
		}
	\end{frame}
	
	% ================================================================
	% Q40: Hash Function in Blockchain
	% ================================================================
	\begin{frame}
		\frametitle{\fontsize{12}{14}\selectfont \textbf{Q40: Hash Function — Blockchain Block Identifier}}
		\begin{block}
			\scriptsize
			What unique function does the hash perform in a blockchain block?
			\vspace{0.1cm}
			\begin{itemize}
				\item[\textbf{A.}] Encrypts all transactions in the block.
				\item[\textbf{B.}] Acts as a unique block identifier.
				\item[\textbf{C.}] Stores the private keys of all users.
				\item[\textbf{D.}] Creates cryptocurrency tokens.
			\end{itemize}
		\end{block}
		\vspace{0.05cm}
		\answerbox{\large \textbf{B} is correct.}
		\vspace{0.05cm}
		\stepbox{\scriptsize
			\keypoint{Explanation:} The hash in a blockchain block \textbf{acts as a unique block identifier}. It is a cryptographic fingerprint of the block's contents. If any data in the block changes, the hash changes, making it tamper-evident. Hashes do NOT encrypt transactions, store private keys, or create tokens.
		}
	\end{frame}
	
	% ================================================================
	% SECTION 23: concentral risk.png
	% ================================================================
	\begin{frame}
		\frametitle{\fontsize{12}{14}\selectfont \textbf{Q41: Concentration Risk — Consequences}}
		\begin{block}
			\scriptsize
			What are common consequences of failing to address concentration risk in financial institutions? (Select Two.)
			\vspace{0.1cm}
			\begin{itemize}
				\item[\textbf{A.}] Regulatory fines for poor credit modeling
				\item[\textbf{B.}] Increased exposure if a single counterparty defaults
				\item[\textbf{C.}] Decreased resilience in times of economic downturn
				\item[\textbf{D.}] Expansion into multiple business lines
			\end{itemize}
		\end{block}
		\vspace{0.05cm}
		\answerbox{\large \textbf{B \& C} are correct.}
		\vspace{0.05cm}
		\stepbox{\scriptsize
			\keypoint{Explanation:}
			\begin{itemize}
				\item \textbf{B:} Concentration risk increases exposure if a single counterparty defaults
				\item \textbf{C:} Decreased resilience in economic downturns (lack of diversification)
			\end{itemize}
			\wronganswer{Fines for poor credit modeling} are a consequence of credit risk, not concentration risk specifically. \wronganswer{Expansion into multiple business lines} is diversification, not a consequence of failing to address concentration risk.
		}
	\end{frame}
	
	% ================================================================
	% SECTION 24: data mining.png
	% ================================================================
	\begin{frame}
		\frametitle{\fontsize{12}{14}\selectfont \textbf{Q42: AFC Data Mining Functions}}
		\begin{block}
			\scriptsize
			What are key functions performed during AFC data mining and matching? (Select Three.)
			\vspace{0.1cm}
			\begin{itemize}
				\item[\textbf{A.}] Uncovering patterns in large complex datasets
				\item[\textbf{B.}] Comparing data from different sources to identify inconsistencies
				\item[\textbf{C.}] Detecting anomalies indicative of fraudulent behavior
				\item[\textbf{D.}] Maintaining databases of known criminal activities
			\end{itemize}
		\end{block}
		\vspace{0.05cm}
		\answerbox{\large \textbf{A, B \& C} are correct.}
		\vspace{0.05cm}
		\stepbox{\scriptsize
			\keypoint{Explanation:}
			\begin{itemize}
				\item \textbf{A:} Pattern discovery in large datasets
				\item \textbf{B:} Cross-source data comparison for inconsistencies
				\item \textbf{C:} Anomaly detection
			\end{itemize}
			\wronganswer{Maintaining databases of criminal activities} is a separate function (watchlist management), not a data mining function.
		}
	\end{frame}
	
	% ================================================================
	% SECTION 25: decline trust.png
	% ================================================================
	\begin{frame}
		\frametitle{\fontsize{12}{14}\selectfont \textbf{Q43: Free Trade Zone — ML Risk Factor}}
		\begin{block}
			\scriptsize
			Which factor most significantly increases the risk of money laundering in free-trade zones (FTZs)?
			\vspace{0.1cm}
			\begin{itemize}
				\item[\textbf{A.}] High corporate tax rates in the FTZ.
				\item[\textbf{B.}] Absence of standardized inspection protocols for transshipped goods.
				\item[\textbf{C.}] Use of a single currency in all transactions.
				\item[\textbf{D.}] Limited access to international shipping.
			\end{itemize}
		\end{block}
		\vspace{0.05cm}
		\answerbox{\large \textbf{B} is correct.}
		\vspace{0.05cm}
		\stepbox{\scriptsize
			\keypoint{Explanation:} The \textbf{absence of standardized inspection protocols for transshipped goods} is the most significant ML risk in FTZs. This allows misclassification, over/under-invoicing, and phantom shipments to go undetected. Tax rates, currency, and shipping access are less relevant to ML risk.
		}
	\end{frame}
	
	% ================================================================
	% Q44: TCSP Client Acceptance
	% ================================================================
	\begin{frame}
		\frametitle{\fontsize{12}{14}\selectfont \textbf{Q44: TCSP Client Acceptance — BO Verification}}
		\begin{block}
			\scriptsize
			A trust and company service provider (TCSP) firm is reviewing whether to accept a client who wants to incorporate in a low-tax Caribbean jurisdiction. The potential client requests confidentiality protections for all shareholders. A connection is identified to industries marked high risk in the NRA. What should the TCSP response be?
			\vspace{0.1cm}
			\begin{itemize}
				\item[\textbf{A.}] Refer to offshore local counsel and detach the firm from compliance responsibility.
				\item[\textbf{B.}] Decline immediately unless full beneficial owner information and risk justifications are disclosed and verified.
				\item[\textbf{C.}] Accept the client with standard CDD.
				\item[\textbf{D.}] Accept the client but monitor for 12 months.
			\end{itemize}
		\end{block}
		\vspace{0.05cm}
		\answerbox{\large \textbf{B} is correct.}
		\vspace{0.05cm}
		\stepbox{\scriptsize
			\keypoint{Explanation:} With confidentiality requests, a low-tax jurisdiction, and connection to high-risk industries, the TCSP must \textbf{decline unless full beneficial owner information and risk justifications are disclosed and verified}. FATF Recs 24/25 require beneficial ownership identification for legal persons and arrangements. Confidentiality cannot override AML obligations.
		}
	\end{frame}
	
	% ================================================================
	% SECTION 26: dnfb.png
	% ================================================================
	\begin{frame}
		\frametitle{\fontsize{12}{14}\selectfont \textbf{Q45: DNFBP Risk Management Strategies}}
		\begin{block}
			\scriptsize
			Which of the following strategies can help institutions manage the risks associated with designated nonfinancial businesses and professions (DNFBPs)? (Select Two.)
			\vspace{0.1cm}
			\begin{itemize}
				\item[\textbf{A.}] Defaulting to enhanced due diligence for all DNFBPs
				\item[\textbf{B.}] Creating a reputational risk or client selection committee
				\item[\textbf{C.}] Using standardized sector-specific questionnaires
				\item[\textbf{D.}] Avoiding all DNFBP relationships
			\end{itemize}
		\end{block}
		\vspace{0.05cm}
		\answerbox{\large \textbf{B \& C} are correct.}
		\vspace{0.05cm}
		\stepbox{\scriptsize
			\keypoint{Explanation:}
			\begin{itemize}
				\item \textbf{Creating a reputational risk or client selection committee} — provides governance for high-risk DNFBP relationships
				\item \textbf{Using standardized sector-specific questionnaires} — enables consistent risk assessment across DNFBPs
			\end{itemize}
			\wronganswer{Defaulting to EDD for all DNFBPs} is inefficient and not risk-based. \wronganswer{Avoiding all DNFBP relationships} is de-risking, which FATF discourages.
		}
	\end{frame}
	
	% ================================================================
	% Q46: Trade Mispricing
	% ================================================================
	\begin{frame}
		\frametitle{\fontsize{12}{14}\selectfont \textbf{Q46: Trade Mispricing — Intentional ML}}
		\begin{block}
			\scriptsize
			A small textile exporter regularly sends shipments priced significantly below market rates, and the foreign counterpart pays the remainder via separate wire transfers to a third party. What does this behavior suggest?
			\vspace{0.1cm}
			\begin{itemize}
				\item[\textbf{A.}] Normal business practice for textile exports.
				\item[\textbf{B.}] Intentional trade mispricing for money laundering.
				\item[\textbf{C.}] Legitimate tax avoidance through transfer pricing.
				\item[\textbf{D.}] Currency hedging strategy.
			\end{itemize}
		\end{block}
		\vspace{0.05cm}
		\answerbox{\large \textbf{B} is correct.}
		\vspace{0.05cm}
		\stepbox{\scriptsize
			\keypoint{Explanation:} Pricing significantly below market + separate wire transfers to a third party = \textbf{intentional trade mispricing for money laundering}. This is a classic TBML technique (under-invoicing) where the value difference is laundered through separate payments. Normal business, tax avoidance, or hedging would not involve third-party payments.
		}
	\end{frame}
	
	% ================================================================
	% SECTION 27: EU AMLA.png
	% ================================================================
	\begin{frame}
		\frametitle{\fontsize{12}{14}\selectfont \textbf{Q47: EU AMLA — Cross-Border Fintech}}
		\begin{block}
			\scriptsize
			An EU-based fintech firm operates online account services in five member states and has been flagged for noncompliance by more than one regulator. Which supervisory authority is most likely to assume oversight?
			\vspace{0.1cm}
			\begin{itemize}
				\item[\textbf{A.}] The national regulator of the member state where the company is headquartered.
				\item[\textbf{B.}] The AMLA, due to the firm's cross-border nature and risk profile.
				\item[\textbf{C.}] The European Central Bank.
				\item[\textbf{D.}] The national regulator of the member state with the most customers.
			\end{itemize}
		\end{block}
		\vspace{0.05cm}
		\answerbox{\large \textbf{B} is correct.}
		\vspace{0.05cm}
		\stepbox{\scriptsize
			\keypoint{Explanation:} The \textbf{AMLA (Anti-Money Laundering Authority)} assumes oversight for cross-border financial institutions with operations in multiple member states, especially when noncompliance has been flagged by more than one regulator. This consolidates supervision and avoids fragmented oversight. National regulators cannot effectively supervise cross-border firms alone.
		}
	\end{frame}
	
	% ================================================================
	% Q48: UK Risk-Based Approach
	% ================================================================
	\begin{frame}
		\frametitle{\fontsize{12}{14}\selectfont \textbf{Q48: UK AML Framework — Risk-Based Approach}}
		\begin{block}
			\scriptsize
			Under the UK's AML framework, what is the primary purpose of the risk-based approach?
			\vspace{0.1cm}
			\begin{itemize}
				\item[\textbf{A.}] To eliminate all AML controls to reduce costs.
				\item[\textbf{B.}] To allow firms to customize controls suited to their risk exposure.
				\item[\textbf{C.}] To require the same controls for all firms regardless of risk.
				\item[\textbf{D.}] To delegate all AML responsibilities to regulators.
			\end{itemize}
		\end{block}
		\vspace{0.05cm}
		\answerbox{\large \textbf{B} is correct.}
		\vspace{0.05cm}
		\stepbox{\scriptsize
			\keypoint{Explanation:} The risk-based approach allows firms to \textbf{customize controls suited to their risk exposure}, allocating resources where they are most needed. It does NOT eliminate controls, require uniform controls, or delegate responsibilities to regulators. The UK's AML framework is based on FATF standards where the risk-based approach is central.
		}
	\end{frame}
	
	% ================================================================
	% SECTION 28: integrating national data.png
	% ================================================================
	\begin{frame}
		\frametitle{\fontsize{12}{14}\selectfont \textbf{Q49: Passive vs Active Liveness Detection}}
		\begin{block}
			\scriptsize
			What distinguishes passive liveness detection from active methods?
			\vspace{0.1cm}
			\begin{itemize}
				\item[\textbf{A.}] Passive liveness requires users to perform specific actions.
				\item[\textbf{B.}] Passive liveness uses minimal user interaction and analyzes subtle facial movement or texture patterns.
				\item[\textbf{C.}] Passive liveness is less secure than active methods.
				\item[\textbf{D.}] Passive liveness requires a physical token.
			\end{itemize}
		\end{block}
		\vspace{0.05cm}
		\answerbox{\large \textbf{B} is correct.}
		\vspace{0.05cm}
		\stepbox{\scriptsize
			\keypoint{Explanation:} \textbf{Passive liveness detection} uses \textbf{minimal user interaction and analyzes subtle facial movement or texture patterns} (micro-expressions, skin texture, eye movement). \textbf{Active methods} require user actions (blinking, turning head, smiling). Passive is more user-friendly but can be less secure against sophisticated spoofing.
		}
	\end{frame}
	
	% ================================================================
	% Q50: Fintech Customer Identity Verification
	% ================================================================
	\begin{frame}
		\frametitle{\fontsize{12}{14}\selectfont \textbf{Q50: Fintech Identity Verification — Additional Measures}}
		\begin{block}
			\scriptsize
			A regulator is assessing a fintech startup's customer identity verification program. The firm currently relies solely on device recognition technology for customer authentication. What additional measures could the fintech implement at initial customer onboarding to improve its financial crime controls? (Select Two.)
			\vspace{0.1cm}
			\begin{itemize}
				\item[\textbf{A.}] Implementing document verification with biometric authentication
				\item[\textbf{B.}] Enhancing digital footprint analysis through IP address verification
				\item[\textbf{C.}] Integrating with national electronic identity verification systems
				\item[\textbf{D.}] Relying solely on social media verification
			\end{itemize}
		\end{block}
		\vspace{0.05cm}
		\answerbox{\large \textbf{A \& C} are correct.}
		\vspace{0.05cm}
		\stepbox{\scriptsize
			\keypoint{Explanation:}
			\begin{itemize}
				\item \textbf{Document verification with biometric authentication} — ensures identity authenticity
				\item \textbf{National electronic identity verification systems} — leverages trusted government identity databases
			\end{itemize}
			\wronganswer{IP address verification} is already part of device recognition. \wronganswer{Social media verification} is not a reliable identity verification method.
		}
	\end{frame}
	
	% ================================================================
	% SECTION 29: international trade topology.png
	% ================================================================
	\begin{frame}
		\frametitle{\fontsize{12}{14}\selectfont \textbf{Q51: TBML — Inaccurate Shipping Documents}}
		\begin{block}
			\scriptsize
			A compliance officer at a multinational bank is asked to brief a new group of employees on how financial crime can manifest in international trade environments. Which of the following best illustrates how financial crime may occur in such contexts?
			\vspace{0.1cm}
			\begin{itemize}
				\item[\textbf{A.}] Trade surplus reports can reveal profitability ratios for various exported goods.
				\item[\textbf{B.}] Shipping documents may contain inaccurate information about the true countries involved.
				\item[\textbf{C.}] Currency exchange rates determine trade volumes.
				\item[\textbf{D.}] Import duties are always paid in cash.
			\end{itemize}
		\end{block}
		\vspace{0.05cm}
		\answerbox{\large \textbf{B} is correct.}
		\vspace{0.05cm}
		\stepbox{\scriptsize
			\keypoint{Explanation:} \textbf{Inaccurate information in shipping documents about true countries involved} is a classic TBML technique. This allows criminals to circumvent sanctions, misrepresent goods, or conceal origin/destination for money laundering. Trade surplus reports, exchange rates, and cash duties are not direct illustrations of financial crime in trade.
		}
	\end{frame}
	
	% ================================================================
	% Q52: Structuring — Small Cash Deposits
	% ================================================================
	\begin{frame}
		\frametitle{\fontsize{12}{14}\selectfont \textbf{Q52: Structuring — Small Cash Deposits}}
		\begin{block}
			\scriptsize
			A fraud analyst discovers a set of small cash deposits made into an account over several days, followed by a sudden international wire transfer. What laundering method is most likely being used?
			\vspace{0.1cm}
			\begin{itemize}
				\item[\textbf{A.}] Trade-based money laundering
				\item[\textbf{B.}] Structuring
				\item[\textbf{C.}] Layering
				\item[\textbf{D.}] Tax evasion
			\end{itemize}
		\end{block}
		\vspace{0.05cm}
		\answerbox{\large \textbf{B} is correct.}
		\vspace{0.05cm}
		\stepbox{\scriptsize
			\keypoint{Explanation:} \textbf{Small cash deposits over several days} (to avoid CTR thresholds) followed by a \textbf{sudden international wire transfer} is classic \textbf{structuring}. The pattern shows deliberate breaking up of cash deposits (smurfing) followed by integration through wire transfer.
		}
	\end{frame}
	
	% ================================================================
	% SECTION 30: layering.png
	% ================================================================
	\begin{frame}
		\frametitle{\fontsize{12}{14}\selectfont \textbf{Q53: Layering — Shell Company Loan}}
		\begin{block}
			\scriptsize
			During a review of a transaction involving a private loan extended to a shell company that has no identifiable operations or revenue, the loan amount matches known criminal proceeds. What is the most appropriate initial classification of this transaction in the context of money laundering?
			\vspace{0.1cm}
			\begin{itemize}
				\item[\textbf{A.}] A placement strategy to integrate clean capital
				\item[\textbf{B.}] A layering tactic designed to conceal the origin of illicit funds
				\item[\textbf{C.}] A legitimate business transaction
				\item[\textbf{D.}] A tax planning strategy
			\end{itemize}
		\end{block}
		\vspace{0.05cm}
		\answerbox{\large \textbf{B} is correct.}
		\vspace{0.05cm}
		\stepbox{\scriptsize
			\keypoint{Explanation:} A \textbf{private loan to a shell company with no operations}, where the loan amount matches criminal proceeds, is a \textbf{layering tactic} to conceal the origin of illicit funds. The loan creates a false paper trail that appears legitimate while distancing funds from their criminal source. This is not placement (which is initial entry into the financial system) or integration (which would involve using the funds legitimately).
		}
	\end{frame}
	
	% ================================================================
	% Q54: Proactive Institutional Response to AML Risks
	% ================================================================
	\begin{frame}
		\frametitle{\fontsize{12}{14}\selectfont \textbf{Q54: Proactive Response to New AML Risks}}
		\begin{block}
			\scriptsize
			What is typically a proactive institutional response to newly detected AML risks?
			\vspace{0.1cm}
			\begin{itemize}
				\item[\textbf{A.}] Reducing reliance on transaction monitoring systems
				\item[\textbf{B.}] Immediately updating the enterprise-wide risk assessment
				\item[\textbf{C.}] Waiting for regulatory guidance before taking action
				\item[\textbf{D.}] Ignoring new risks until they materialize
			\end{itemize}
		\end{block}
		\vspace{0.05cm}
		\answerbox{\large \textbf{B} is correct.}
		\vspace{0.05cm}
		\stepbox{\scriptsize
			\keypoint{Explanation:} A proactive response is to \textbf{immediately update the enterprise-wide risk assessment} when new AML risks are detected. This allows the institution to adjust controls, allocate resources, and enhance monitoring before risks materialize. Reducing reliance on monitoring, waiting for guidance, or ignoring risks are reactive and not proactive.
		}
	\end{frame}
	
	% ================================================================
	% SECTION 31: oafc.png
	% ================================================================
	\begin{frame}
		\frametitle{\fontsize{12}{14}\selectfont \textbf{Q55: Rules-Based Monitoring — Legacy Customer}}
		\begin{block}
			\scriptsize
			A rules-based transaction monitoring system generates a high volume of alerts based on large deposits. However, the compliance team discovers that a significant number of these alerts relate to a single legacy corporate customer whose monthly deposits have remained consistent for years. What is the most appropriate next step?
			\vspace{0.1cm}
			\begin{itemize}
				\item[\textbf{A.}] Escalate all alerts to the Level 2 team for manual investigation.
				\item[\textbf{B.}] Update customer information about expected transaction patterns in the monitoring system.
				\item[\textbf{C.}] Ignore the alerts as they are false positives.
				\item[\textbf{D.}] Increase the alert threshold for this customer.
			\end{itemize}
		\end{block}
		\vspace{0.05cm}
		\answerbox{\large \textbf{B} is correct.}
		\vspace{0.05cm}
		\stepbox{\scriptsize
			\keypoint{Explanation:} When a legacy customer with consistent historical activity generates alerts, the most appropriate step is to \textbf{update customer information about expected transaction patterns in the monitoring system}. This reduces false positives while maintaining detection for genuinely suspicious activity. Escalating all alerts is inefficient, ignoring is negligent, and increasing thresholds may miss real risk.
		}
	\end{frame}
	
	% ================================================================
	% Q56: Sanctions Screening for Small Business
	% ================================================================
	\begin{frame}
		\frametitle{\fontsize{12}{14}\selectfont \textbf{Q56: Sanctions Screening — Small Business Setup}}
		\begin{block}
			\scriptsize
			A luxury brand based in the United States operates online with a modest client base and limited financial volume. They want to implement sanctions screening. Which setup is most appropriate given their profile?
			\vspace{0.1cm}
			\begin{itemize}
				\item[\textbf{A.}] Implementing a batch screening audit conducted annually by external consultants
				\item[\textbf{B.}] Using spreadsheets linked to the OFAC website for customer checks
				\item[\textbf{C.}] Implementing a real-time automated screening system with full integration
				\item[\textbf{D.}] Relying on the payment processor for all screening
			\end{itemize}
		\end{block}
		\vspace{0.05cm}
		\answerbox{\large \textbf{B} is correct.}
		\vspace{0.05cm}
		\stepbox{\scriptsize
			\keypoint{Explanation:} For a small business with modest volume, \textbf{using spreadsheets linked to the OFAC website for customer checks} is proportionate and appropriate. The risk-based approach allows scaled controls based on size and risk profile. Full automated systems may be cost-prohibitive, and relying solely on payment processors may not cover all AML obligations.
		}
	\end{frame}
	
	% ================================================================
	% SECTION 32: biometrics.png
	% ================================================================
	\begin{frame}
		\frametitle{\fontsize{12}{14}\selectfont \textbf{Q57: Facial Recognition — Compromised Data Risk}}
		\begin{block}
			\scriptsize
			Which risk is specifically associated with facial recognition technology compared to traditional password-based authentication methods?
			\vspace{0.1cm}
			\begin{itemize}
				\item[\textbf{A.}] The technology requires more storage capacity than traditional methods.
				\item[\textbf{B.}] Compromised biometric data is difficult to change or reset unlike passwords.
				\item[\textbf{C.}] Facial recognition is less secure than passwords.
				\item[\textbf{D.}] Facial recognition requires complex hardware.
			\end{itemize}
		\end{block}
		\vspace{0.05cm}
		\answerbox{\large \textbf{B} is correct.}
		\vspace{0.05cm}
		\stepbox{\scriptsize
			\keypoint{Explanation:} The key risk with facial recognition is that \textbf{compromised biometric data is difficult to change or reset}, unlike passwords which can be changed. Once a biometric template is compromised, it cannot be replaced. Storage capacity, hardware complexity, and security compared to passwords are not unique risks of facial recognition.
		}
	\end{frame}
	
	% ================================================================
	% Q58: Authentication Technologies
	% ================================================================
	\begin{frame}
		\frametitle{\fontsize{12}{14}\selectfont \textbf{Q58: Authentication Technologies}}
		\begin{block}
			\scriptsize
			Which technologies enhance authentication and security in financial systems? (Select Three.)
			\vspace{0.1cm}
			\begin{itemize}
				\item[\textbf{A.}] Multi-factor authentication
				\item[\textbf{B.}] Physiological biometrics
				\item[\textbf{C.}] Behavioral biometrics
				\item[\textbf{D.}] Password biometrics
			\end{itemize}
		\end{block}
		\vspace{0.05cm}
		\answerbox{\large \textbf{A, B \& C} are correct.}
		\vspace{0.05cm}
		\stepbox{\scriptsize
			\keypoint{Explanation:}
			\begin{itemize}
				\item \textbf{Multi-factor authentication} — combines multiple authentication factors
				\item \textbf{Physiological biometrics} — physical traits (fingerprints, facial features, iris)
				\item \textbf{Behavioral biometrics} — behavioral patterns (typing rhythm, mouse movements)
			\end{itemize}
			\wronganswer{Password biometrics} is not a recognized authentication category — passwords are not biometrics.
		}
	\end{frame}
	
	% ================================================================
	% SECTION 33: board members.png
	% ================================================================
	\begin{frame}
		\frametitle{\fontsize{12}{14}\selectfont \textbf{Q59: Board Oversight — Compliance Culture}}
		\begin{block}
			\scriptsize
			An internal compliance audit finds that despite robust policies, employee engagement on compliance is low. Survey data shows most staff view compliance as a barrier to performance. Which conclusion is most accurate?
			\vspace{0.1cm}
			\begin{itemize}
				\item[\textbf{A.}] Compliance policies are too complex for employees to understand.
				\item[\textbf{B.}] The organization has failed to embed compliance values in its operational culture.
				\item[\textbf{C.}] Employees need more training on compliance policies.
				\item[\textbf{D.}] The audit findings are not significant.
			\end{itemize}
		\end{block}
		\vspace{0.05cm}
		\answerbox{\large \textbf{B} is correct.}
		\vspace{0.05cm}
		\stepbox{\scriptsize
			\keypoint{Explanation:} Low employee engagement despite robust policies indicates the organization has \textbf{failed to embed compliance values in its operational culture}. Policies alone are insufficient — culture drives behavior. While training or simplification might help, the root issue is cultural, not just informational.
		}
	\end{frame}
	
	% ================================================================
	% Q60: Board Oversight — AFC Program
	% ================================================================
	\begin{frame}
		\frametitle{\fontsize{12}{14}\selectfont \textbf{Q60: Board Oversight of AFC Program}}
		\begin{block}
			\scriptsize
			Which of the following are ways board members demonstrate effective oversight over an institution's AFC program? (Select Three.)
			\vspace{0.1cm}
			\begin{itemize}
				\item[\textbf{A.}] Participating in SAR investigations to understand typologies
				\item[\textbf{B.}] Receiving regular reports on financial crime risks and control effectiveness
				\item[\textbf{C.}] Reviewing whether senior management is addressing audit findings
				\item[\textbf{D.}] Challenging management when AFC performance metrics show signs of weakness
			\end{itemize}
		\end{block}
		\vspace{0.05cm}
		\answerbox{\large \textbf{B, C \& D} are correct.}
		\vspace{0.05cm}
		\stepbox{\scriptsize
			\keypoint{Explanation:}
			\begin{itemize}
				\item \textbf{B:} Receiving regular reports on risks and control effectiveness
				\item \textbf{C:} Reviewing management's response to audit findings
				\item \textbf{D:} Challenging management on performance metrics
			\end{itemize}
			\wronganswer{Participating in SAR investigations} is operational and not a board-level function — boards provide oversight, not operational involvement.
		}
	\end{frame}
	
	% ================================================================
	% SECTION 34: concentration account.png
	% ================================================================
	\begin{frame}
		\frametitle{\fontsize{12}{14}\selectfont \textbf{Q61: Concentration Account — Regulatory Concern}}
		\begin{block}
			\scriptsize
			During a routine compliance review, a financial institution finds that its treasury unit uses a concentration account that lacks sufficient oversight. What is the most likely regulatory concern?
			\vspace{0.1cm}
			\begin{itemize}
				\item[\textbf{A.}] The account has high transaction fees.
				\item[\textbf{B.}] The institution might inadvertently facilitate money laundering due to poor traceability.
				\item[\textbf{C.}] The account violates tax laws.
				\item[\textbf{D.}] The account exceeds interest rate limits.
			\end{itemize}
		\end{block}
		\vspace{0.05cm}
		\answerbox{\large \textbf{B} is correct.}
		\vspace{0.05cm}
		\stepbox{\scriptsize
			\keypoint{Explanation:} Concentration accounts (nostro/vostro accounts) with insufficient oversight create \textbf{poor traceability of funds}, making it difficult to identify the ultimate originator or beneficiary of transactions. This \textbf{inadvertently facilitates money laundering} by obscuring the source of funds. Fee rates, tax issues, and interest rates are not the primary regulatory concern.
		}
	\end{frame}
	
	% ================================================================
	% SECTION 35: national electronic verification system.png
	% ================================================================
	\begin{frame}
		\frametitle{\fontsize{12}{14}\selectfont \textbf{Q62: National e-ID Integration}}
		\begin{block}
			\scriptsize
			A regulator is assessing a fintech startup's customer identity verification program. The firm currently relies solely on device recognition technology for customer authentication. What additional measures could the fintech implement at initial customer onboarding to improve its financial crime controls? (Select Two.)
			\vspace{0.1cm}
			\begin{itemize}
				\item[\textbf{A.}] Implementing document verification with biometric authentication
				\item[\textbf{B.}] Integrating with national electronic identity verification systems
				\item[\textbf{C.}] Enhancing digital footprint analysis through IP address verification
				\item[\textbf{D.}] Relying solely on credit bureau data
			\end{itemize}
		\end{block}
		\vspace{0.05cm}
		\answerbox{\large \textbf{A \& B} are correct.}
		\vspace{0.05cm}
		\stepbox{\scriptsize
			\keypoint{Explanation:}
			\begin{itemize}
				\item \textbf{Document verification with biometric authentication} — ensures identity authenticity
				\item \textbf{National electronic identity verification systems} — leverages government identity databases
			\end{itemize}
			IP verification is already part of device recognition, and credit bureau data is not a primary identity verification method.
		}
	\end{frame}
	
	% ================================================================
	% SECTION 36: terrori humant traffic topology.png
	% ================================================================
	\begin{frame}
		\frametitle{\fontsize{12}{14}\selectfont \textbf{Q63: Terrorist Financing — Movement Methods}}
		\begin{block}
			\scriptsize
			Which of the following techniques are commonly used by terrorists to move or store money across borders while avoiding detection? (Select Three.)
			\vspace{0.1cm}
			\begin{itemize}
				\item[\textbf{A.}] Establishing fake charities in high-risk regions
				\item[\textbf{B.}] Use of informal value transfer systems like Hawala
				\item[\textbf{C.}] Physical transportation of cash across porous borders
				\item[\textbf{D.}] Large wire transfers through regulated banks
			\end{itemize}
		\end{block}
		\vspace{0.05cm}
		\answerbox{\large \textbf{A, B \& C} are correct.}
		\vspace{0.05cm}
		\stepbox{\scriptsize
			\keypoint{Explanation:}
			\begin{itemize}
				\item \textbf{Fake charities} — disguise terrorist financing as humanitarian aid
				\item \textbf{Hawala} — informal value transfer outside regulated banking
				\item \textbf{Physical cash transport} — crosses borders avoiding detection
			\end{itemize}
			\wronganswer{Large wire transfers through regulated banks} are more easily detected and monitored, making them less attractive for terrorists.
		}
	\end{frame}
	
	% ================================================================
	% Q64: Human Trafficking — Laundering Proceeds
	% ================================================================
	\begin{frame}
		\frametitle{\fontsize{12}{14}\selectfont \textbf{Q64: Human Trafficking — Laundering Illicit Proceeds}}
		\begin{block}
			\scriptsize
			Which of the following characteristics enable human trafficking operations to effectively launder illicit proceeds? (Select Two.)
			\vspace{0.1cm}
			\begin{itemize}
				\item[\textbf{A.}] High correlation between trafficking and counterfeit document operations
				\item[\textbf{B.}] Use of mobile payment and remittance systems
				\item[\textbf{C.}] Weak cross-border enforcement cooperation
				\item[\textbf{D.}] Consistent global anti-trafficking enforcement
			\end{itemize}
		\end{block}
		\vspace{0.05cm}
		\answerbox{\large \textbf{B \& C} are correct.}
		\vspace{0.05cm}
		\stepbox{\scriptsize
			\keypoint{Explanation:}
			\begin{itemize}
				\item \textbf{Mobile payment and remittance systems} — enable rapid movement of small amounts (structuring)
				\item \textbf{Weak cross-border enforcement cooperation} — allows traffickers to operate across jurisdictions
			\end{itemize}
			\wronganswer{Counterfeit documents} are a separate crime, not specifically enabling laundering. \wronganswer{Consistent enforcement} would reduce, not enable, laundering.
		}
	\end{frame}
	
	% ================================================================
	% SECTION 37: trade topology.png
	% ================================================================
	\begin{frame}
		\frametitle{\fontsize{12}{14}\selectfont \textbf{Q65: Penny Stock Manipulation Risk}}
		\begin{block}
			\scriptsize
			A compliance officer is evaluating two trades: one involves high-volume equities with normal volatility; the other, low-volume penny stocks with rapid pricing shifts. Both trades passed basic monitoring. Which should be prioritized for additional review?
			\vspace{0.1cm}
			\begin{itemize}
				\item[\textbf{A.}] Both equally, as volume and volatility are neutral indicators.
				\item[\textbf{B.}] The penny stock trade, due to higher manipulation risk.
				\item[\textbf{C.}] The high-volume equity trade, due to larger potential impact.
				\item[\textbf{D.}] Neither, as both passed basic monitoring.
			\end{itemize}
		\end{block}
		\vspace{0.05cm}
		\answerbox{\large \textbf{B} is correct.}
		\vspace{0.05cm}
		\stepbox{\scriptsize
			\keypoint{Explanation:} \textbf{Low-volume penny stocks with rapid pricing shifts} present higher manipulation risk (pump-and-dump, wash trading). High-volume equities with normal volatility are less suspicious. Passing basic monitoring does not eliminate the need for risk-based prioritization.
		}
	\end{frame}
	
	% ================================================================
	% SECTION 38: batch screening.png
	% ================================================================
	\begin{frame}
		\frametitle{\fontsize{12}{14}\selectfont \textbf{Q66: Batch Screening — Benefits}}
		\begin{block}
			\scriptsize
			Which of the following are true about batch screening? (Select Two.)
			\vspace{0.1cm}
			\begin{itemize}
				\item[\textbf{A.}] It helps reveal changing risks within the existing customer base.
				\item[\textbf{B.}] It is designed to handle large volumes of data.
				\item[\textbf{C.}] It provides real-time transaction screening.
				\item[\textbf{D.}] It eliminates the need for ongoing monitoring.
			\end{itemize}
		\end{block}
		\vspace{0.05cm}
		\answerbox{\large \textbf{A \& B} are correct.}
		\vspace{0.05cm}
		\stepbox{\scriptsize
			\keypoint{Explanation:}
			\begin{itemize}
				\item \textbf{Reveals changing risks within existing customer base} — periodic re-screening catches new designations
				\item \textbf{Handles large volumes of data} — efficient for large customer databases
			\end{itemize}
			\wronganswer{Real-time screening} is not batch screening. \wronganswer{Eliminates ongoing monitoring} is false — batch screening supplements ongoing monitoring.
		}
	\end{frame}
	
	% ================================================================
	% SECTION 39: cust onboarding.png
	% ================================================================
	\begin{frame}
		\frametitle{\fontsize{12}{14}\selectfont \textbf{Q67: Risk-Based Authentication — User Experience}}
		\begin{block}
			\scriptsize
			How do risk-based authentication systems optimize user experience?
			\vspace{0.1cm}
			\begin{itemize}
				\item[\textbf{A.}] They eliminate all authentication requirements.
				\item[\textbf{B.}] They prompt additional checks only when user behavior deviates from expectations.
				\item[\textbf{C.}] They always require multi-factor authentication for all users.
				\item[\textbf{D.}] They allow all users to bypass authentication.
			\end{itemize}
		\end{block}
		\vspace{0.05cm}
		\answerbox{\large \textbf{B} is correct.}
		\vspace{0.05cm}
		\stepbox{\scriptsize
			\keypoint{Explanation:} Risk-based authentication optimizes user experience by \textbf{prompting additional checks only when user behavior deviates from expectations}. Low-risk users have frictionless experiences, while higher-risk activities trigger additional authentication. This balances security and convenience.
		}
	\end{frame}
	
	% ================================================================
	% SECTION 40: fase positive.png
	% ================================================================
	\begin{frame}
		\frametitle{\fontsize{12}{14}\selectfont \textbf{Q68: Data Extraction Challenges}}
		\begin{block}
			\scriptsize
			Which of the following best describes a challenge during AFC data extraction?
			\vspace{0.1cm}
			\begin{itemize}
				\item[\textbf{A.}] Data always comes in standardized formats.
				\item[\textbf{B.}] Data may come in different formats, requiring normalization during integration.
				\item[\textbf{C.}] Data extraction never affects analysis quality.
				\item[\textbf{D.}] All data sources are perfectly compatible.
			\end{itemize}
		\end{block}
		\vspace{0.05cm}
		\answerbox{\large \textbf{B} is correct.}
		\vspace{0.05cm}
		\stepbox{\scriptsize
			\keypoint{Explanation:} A key challenge in AFC data extraction is that \textbf{data may come in different formats, requiring normalization during integration}. Different systems, jurisdictions, and data standards create compatibility issues. Standardized formats are rare, and compatibility cannot be assumed.
		}
	\end{frame}
	
	% ================================================================
	% Q69: Data Mining Requirements
	% ================================================================
	\begin{frame}
		\frametitle{\fontsize{12}{14}\selectfont \textbf{Q69: Data Mining Requirements}}
		\begin{block}
			\scriptsize
			Which of the following are requirements for effective data mining and matching in AFC controls? (Select Two.)
			\vspace{0.1cm}
			\begin{itemize}
				\item[\textbf{A.}] Comparing extracted data against structured and unstructured risk sources
				\item[\textbf{B.}] Focusing on internal customer files for mining purposes
				\item[\textbf{C.}] Applying intelligent filters to reduce false positives during entity resolution
				\item[\textbf{D.}] Excluding all external data sources
			\end{itemize}
		\end{block}
		\vspace{0.05cm}
		\answerbox{\large \textbf{A \& C} are correct.}
		\vspace{0.05cm}
		\stepbox{\scriptsize
			\keypoint{Explanation:}
			\begin{itemize}
				\item \textbf{Comparing against structured and unstructured risk sources} — broadens detection
				\item \textbf{Intelligent filters to reduce false positives} — improves efficiency
			\end{itemize}
			\wronganswer{Focusing only on internal files} and \wronganswer{excluding external sources} would undermine data mining effectiveness.
		}
	\end{frame}
	
	% ================================================================
	% SECTION 41: mlat.png
	% ================================================================
	\begin{frame}
		\frametitle{\fontsize{12}{14}\selectfont \textbf{Q70: MLAT Implementation Challenges}}
		\begin{block}
			\scriptsize
			Two countries, A and B, are attempting to establish a mutual legal assistance treaty. Country A has data protection rules preventing sharing of client identities, while Country B bases most financial investigations on beneficial ownership disclosures. What challenge is most likely to arise in implementing the MLAT?
			\vspace{0.1cm}
			\begin{itemize}
				\item[\textbf{A.}] Evidence exchanged under the MLAT will not be admissible in Country B's legal system.
				\item[\textbf{B.}] Country A's data protection rules will prevent sharing client identities.
				\item[\textbf{C.}] Country B will not accept beneficial ownership information.
				\item[\textbf{D.}] The MLAT will be easy to implement.
			\end{itemize}
		\end{block}
		\vspace{0.05cm}
		\answerbox{\large \textbf{B} is correct.}
		\vspace{0.05cm}
		\stepbox{\scriptsize
			\keypoint{Explanation:} \textbf{Country A's data protection rules preventing sharing of client identities} creates a fundamental challenge for MLAT implementation. While Country B requires beneficial ownership information, Country A cannot share it. This conflict between privacy laws and mutual legal assistance is a common MLAT challenge. Admissibility is not the primary issue — the primary issue is data availability.
		}
	\end{frame}
	
	% ================================================================
	% Q71: PPP — Information Sharing
	% ================================================================
	\begin{frame}
		\frametitle{\fontsize{12}{14}\selectfont \textbf{Q71: PPP — Collaboration Model}}
		\begin{block}
			\scriptsize
			Which of the following best describes the principle behind public-private partnerships (PPPs) in AML efforts?
			\vspace{0.1cm}
			\begin{itemize}
				\item[\textbf{A.}] Government agencies share all intelligence with the public.
				\item[\textbf{B.}] Collaboration allows timely information sharing and joint targeting of high-risk actors using shared intelligence.
				\item[\textbf{C.}] Private sector replaces all government enforcement.
				\item[\textbf{D.}] PPPs eliminate the need for regulatory reporting.
			\end{itemize}
		\end{block}
		\vspace{0.05cm}
		\answerbox{\large \textbf{B} is correct.}
		\vspace{0.05cm}
		\stepbox{\scriptsize
			\keypoint{Explanation:} PPPs enable \textbf{collaboration for timely information sharing and joint targeting of high-risk actors}. They supplement, not replace, government enforcement or regulatory reporting. PPPs facilitate information flow between public and private sectors to combat financial crime more effectively.
		}
	\end{frame}
	
	% ================================================================
	% SECTION 42: multi data sources.png
	% ================================================================
	\begin{frame}
		\frametitle{\fontsize{12}{14}\selectfont \textbf{Q72: Multi-Source Data Integration}}
		\begin{block}
			\scriptsize
			Which of the following help ensure effective integration of multiple data sources into compliance platforms? (Select Three.)
			\vspace{0.1cm}
			\begin{itemize}
				\item[\textbf{A.}] Mapping all incoming data to standardized field names
				\item[\textbf{B.}] Maintaining a common data dictionary across systems
				\item[\textbf{C.}] Aligning control data requirements to system inputs
				\item[\textbf{D.}] Separating internal and external data into isolated systems
			\end{itemize}
		\end{block}
		\vspace{0.05cm}
		\answerbox{\large \textbf{A, B \& C} are correct.}
		\vspace{0.05cm}
		\stepbox{\scriptsize
			\keypoint{Explanation:}
			\begin{itemize}
				\item \textbf{Mapping data to standardized field names} — ensures consistency
				\item \textbf{Maintaining a common data dictionary} — provides a single source of truth
				\item \textbf{Aligning control requirements to inputs} — ensures systems capture needed data
			\end{itemize}
			\wronganswer{Separating internal and external data into isolated systems} undermines integration rather than enabling it.
		}
	\end{frame}
	
	% ================================================================
	% Q73: Internal vs External Sources
	% ================================================================
	\begin{frame}
		\frametitle{\fontsize{12}{14}\selectfont \textbf{Q73: Internal vs External Data Sources}}
		\begin{block}
			\scriptsize
			Which characteristics describe internal versus external sources used during AFC data extraction? (Select Two.)
			\vspace{0.1cm}
			\begin{itemize}
				\item[\textbf{A.}] External sources include sanction or adverse media lists from agencies.
				\item[\textbf{B.}] Internal data includes account balance snapshots, not user activity.
				\item[\textbf{C.}] Internal sources reflect observed customer behavior or system-held values.
				\item[\textbf{D.}] External sources include internal transaction histories.
			\end{itemize}
		\end{block}
		\vspace{0.05cm}
		\answerbox{\large \textbf{A \& C} are correct.}
		\vspace{0.05cm}
		\stepbox{\scriptsize
			\keypoint{Explanation:}
			\begin{itemize}
				\item \textbf{External sources:} Sanction lists, adverse media, PEP lists, watchlists
				\item \textbf{Internal sources:} Customer behavior, system-held values, transaction history, account data
			\end{itemize}
			\wronganswer{Account balance snapshots only} is too narrow. \wronganswer{External includes internal histories} is incorrect.
		}
	\end{frame}
	
	% ================================================================
	% SECTION 43: ppp model.png
	% ================================================================
	\begin{frame}
		\frametitle{\fontsize{12}{14}\selectfont \textbf{Q74: PPP — Two-Way Communication}}
		\begin{block}
			\scriptsize
			A national financial intelligence unit (FIU) wants to collaborate with financial institutions to disrupt a growing terrorist financing network. It intends to use the public-private partnership (PPP) model for real-time intelligence sharing. Which of the following approaches would be the most aligned with the goals of a PPP framework?
			\vspace{0.1cm}
			\begin{itemize}
				\item[\textbf{A.}] Designing a protocol in which information is transmitted only in anonymized, aggregated format after a threat is resolved.
				\item[\textbf{B.}] Establishing a formal platform where regulators and institutions engage in direct two-way communication on specific threats.
				\item[\textbf{C.}] Disseminating public reports without engaging financial institutions.
				\item[\textbf{D.}] Requiring institutions to submit data without receiving feedback.
			\end{itemize}
		\end{block}
		\vspace{0.05cm}
		\answerbox{\large \textbf{B} is correct.}
		\vspace{0.05cm}
		\stepbox{\scriptsize
			\keypoint{Explanation:} PPPs require \textbf{direct two-way communication on specific threats} between regulators/FIUs and financial institutions. Anonymized post-threat reports (A), one-way public reports (C), and one-way data submission (D) do NOT provide the real-time collaboration that defines PPPs.
		}
	\end{frame}
	
	% ================================================================
	% Q75: Cross-Border Regulatory Coordination
	% ================================================================
	\begin{frame}
		\frametitle{\fontsize{12}{14}\selectfont \textbf{Q75: Cross-Border Regulatory Coordination}}
		\begin{block}
			\scriptsize
			A cross-border financial institution is supervised by three different regulators across two regions. To maintain comprehensive compliance while reducing administrative burden, what integrated strategy should the regulators apply?
			\vspace{0.1cm}
			\begin{itemize}
				\item[\textbf{A.}] Each regulator conducts independent reviews without coordination.
				\item[\textbf{B.}] Regulators should coordinate to clarify scopes, align risk approaches, and execute joint reviews.
				\item[\textbf{C.}] One regulator takes sole responsibility for all supervision.
				\item[\textbf{D.}] The institution chooses which regulator to report to.
			\end{itemize}
		\end{block}
		\vspace{0.05cm}
		\answerbox{\large \textbf{B} is correct.}
		\vspace{0.05cm}
		\stepbox{\scriptsize
			\keypoint{Explanation:} For cross-border institutions, regulators should \textbf{coordinate to clarify scopes, align risk approaches, and execute joint reviews}. This reduces administrative burden while ensuring comprehensive compliance. Independent reviews are duplicative, sole responsibility may be legally impossible, and allowing the institution to choose violates regulatory authority.
		}
	\end{frame}
	
	% ================================================================
	% EXAM TIPS SLIDE
	% ================================================================
\begin{frame}
	\frametitle{\textbf{Exam Tips — Common CAMS Traps}}
	
	\small
	\begin{itemize}
		\item \textbf{TF vs ML}: TF = intent/destination (can be legit funds). ML = source (always illicit).
		\item \textbf{SAR vs STR}: US = SARs. Non-US DNFBPs = STRs to FIU.
		\item \textbf{Rec 15 vs 16}: 15 = VASP rules. 16 = Travel Rule.
		\item \textbf{Domestic PEPs}: Risk-based; apply EDD if high risk.
		\item \textbf{50\% Rule}: 50\%+ ownership → entity is blocked.
		\item \textbf{Structuring}: Illegal even with clean money.
		\item \textbf{Tipping Off}: Never disclose SAR.
		\item \textbf{SDD}: Not zero CDD; suspicion overrides.
		\item \textbf{Shell Banks}: Always prohibited.
		\item \textbf{De-Risking}: No blanket exits; manage risk.
		\item \textbf{SOW vs SOF}: Different; both required.
		\item \textbf{PF vs TF}: PF = WMD (Rec 7), TF = terrorism (Rec 5).
		\item \textbf{Device Data}: IP, geo, device — not social media.
		\item \textbf{ESG}: Reporting concept, not AML.
		\item \textbf{TBML}: Over = seller benefit; Under = buyer benefit.
	\end{itemize}
	
	\vspace{0.2cm}
	\textbf{Tip:} Ask: \textit{Always, sometimes, or never true?}
\end{frame}
	
	% ================================================================
	% SUMMARY TABLE
	% ================================================================
	\begin{frame}
		\frametitle{\fontsize{12}{14}\selectfont \textbf{Summary: Key Rules for CAMS}}
		\vspace{0.1cm}
		\scriptsize
		\begin{longtable}{|p{0.28\textwidth}|p{0.67\textwidth}|}
			\hline
			\keypoint{Rule} & \keypoint{What the CAMS Exam Tests} \\
			\hline
			SAR vs STR & US BSA: defined FIs file SARs. US lawyers/real estate agents = NO BSA SAR. \\
			\hline
			TF vs ML & TF funds can be legitimate; ML funds always illicit. \\
			\hline
			Sanctions & Strict liability (civil). OFAC 50\% rule. \\
			\hline
			PEPs & Foreign: mandatory EDD. Domestic: risk-based assessment. SOW and SOF both required. \\
			\hline
			VASPs & Foundational = Rec 15. Travel Rule = Rec 16. \\
			\hline
			Structuring & Illegal even with clean money. File SAR; do NOT tip off. \\
			\hline
			Correspondent Banking & EDD + senior approval. Shell banks: absolute prohibition. \\
			\hline
			Beneficial Ownership & Natural person who ultimately owns/controls. 25\% = starting point. \\
			\hline
			FATF Evaluations & Technical compliance + effectiveness. Gray List = EDD (Rec 19). \\
			\hline
			Device Intelligence & IP/geolocation, OS/plugins, browser/screen — NOT social media. \\
			\hline
			TBML & Over-invoicing = value to seller. Under-invoicing = value to buyer. \\
			\hline
			IOSCO & Core objectives: investor protection + financial stability. \\
			\hline
			PPP & Two-way communication between regulators and FIs. \\
			\hline
		\end{longtable}
	\end{frame}
	
	% ================================================================
	% END SLIDE
	% ================================================================
	\begin{frame}
		\centering
		\vspace{1.5cm}
		{\fontsize{18}{22}\selectfont \textbf{Domain A — All PNG Questions}}\\[0.3cm]
		{\fontsize{12}{14}\selectfont All answers verified against FATF Recommendations, US BSA, and CAMS Study Guide.}\\[0.3cm]
		\rule{4cm}{0.5pt}\\[0.2cm]
		{\fontsize{9}{11}\selectfont \textit{75+ questions from all 43 image files — fully explained with correct answers,}\\
			\textit{wrong answer explanations, and exam traps.}}
		\vspace{0.5cm}
	\end{frame}
	
	% ================================================================
	% END DOCUMENT
	% ================================================================
\end{document}